\documentclass[11pt,a4paper]{article}
\usepackage[T1]{fontenc}
\usepackage[utf8]{inputenc}
\usepackage{lmodern,microtype,amsmath,amssymb,amsthm,mathtools,booktabs,array,xcolor}
\usepackage[left=25mm,right=25mm,top=23mm,bottom=25mm]{geometry}
\usepackage[hypertexnames=false]{hyperref}
\definecolor{inkblue}{HTML}{214D68}
\hypersetup{colorlinks=true,linkcolor=inkblue,citecolor=inkblue,urlcolor=inkblue,
pdftitle={Restricted partitions, domination polynomials, and the Kanter factor},pdfauthor={Research manuscript prepared with ChatGPT for Keio}}
\newtheorem{theorem}{Theorem}[section]
\newtheorem{lemma}[theorem]{Lemma}
\newtheorem{proposition}[theorem]{Proposition}
\newtheorem{corollary}[theorem]{Corollary}
\numberwithin{equation}{section}
\setlength{\emergencystretch}{2em}
\setlength{\parskip}{2pt plus 0.5pt}
\title{Restricted partitions, domination polynomials,\protect\\and the Kanter factor}
\author{Research manuscript prepared with ChatGPT for Keio}
\date{3 November 2026}
\begin{document}
\maketitle
\begin{abstract}
Let $p_A(n)$ count partitions of $n$ into parts from a fixed set $A$ of
positive integers with greatest common divisor one. We prove that
$p_A(h+n)/(p_A(h)p_A(n))$ tends to zero uniformly for $n\ge h\to\infty$
whenever $|A|\ge2$. If $A$ is infinite, the ratio decays faster than
every fixed power of $1/h$. The proof cuts a partition at a
chosen sum and bounds the number of partitions that produce the same
two lists. A separate finite calculation determines all exceptions in
the square-partition inequality.

Random deletion from a dominating set proves ultra-log-concavity
for powers of paths and cycles. For trees, removing one leaf can move
the most frequent size of a dominating set by an unbounded distance,
even when both sequences are unimodal. In probability, a contour estimate
and a compact interval certificate give the generalized gamma convolution
range $0<\beta\le4/5$ for the Kanter factor $K_\beta$.
A different prediction for the shifted positive free stable law fails
throughout $23/25\le\beta<1$, where its density cannot be a gamma
scale mixture of shape $3/2$. A shifted gamma integral establishes
the coefficient identity at the limiting stable edge in all orders. Finite bell shape also distinguishes
interior derivative zero counts from endpoint conditions.
The remaining results concern sharp Jacobi endpoint comparison,
global Bessel bounds, and the fractional completely monotonic degrees
of two trigamma remainders.
\end{abstract}
\noindent\textbf{Scope.} Page~\pageref{sec:map} identifies the ten
source statements. Six are proved and four are disproved under their
printed hypotheses. Established special cases and closely related
earlier results are credited where used. The manuscript includes its
five finite verifiers. These are mathematical claims submitted for
scrutiny. Worldwide priority and comparative rank within the entire
source corpus have not been established.
\clearpage
\section{Original questions and proof locations}\label{sec:map}
Entries 3 and 4 involve ordinary domination on two different graph
families. Entries 5--7 ask about two probability laws and an endpoint
expansion in the same source.
A single entry is used for each source conjecture with several clauses.
\begin{center}\small
\begin{tabular}{@{}r>{\raggedright\arraybackslash}p{.318\textwidth}>{\raggedright\arraybackslash}p{.53\textwidth}r@{}}
\toprule
& Source & Conclusion & Page\\
\midrule
1 & Gajdzica, Conj.~6.5 (2024 preprint) & True for arbitrary allowed parts. The ratio decays superpolynomially if the set of parts is infinite. & \pageref{sec:arbitrary-A}\\[4pt]
2 & Roy, Conj.~3.2 (2024 preprint) & True. All equalities and failures for partitions into squares agree with the proposed table. & \pageref{sec:roy-squares}\\[4pt]
3 & Galvin--Zhang, Conj.~1.8 (2024) & True. Domination coefficients divided by the corresponding binomial coefficients are log-concave for every power of both paths and cycles. & \pageref{sec:gz-ulc}\\[4pt]
4 & S.~Zhang, Conj.~4.1 (2021) & False. Removing one leaf can move every mode arbitrarily far while preserving unimodality. & \pageref{sec:tree-leaf-mode}\\[4pt]
5 & Hasebe--Simon--Wang, Conj.~1 (2018/2020) & True. The Kanter variable is a generalized gamma convolution exactly for $\beta\le4/5$. & \pageref{sec:kanter}\\[4pt]
6 & Hasebe--Simon--Wang, Remarks~2(b), 10(b) & False. The shifted positive free stable law has the proposed gamma mixture for $\beta\le4/5$ and fails it throughout $23/25\le\beta<1$. & \pageref{sec:shifted-stable}\\[4pt]
7 & Hasebe--Simon--Wang, Remark~12 & True. The stable-edge coefficient equals the proposed partial sum of Stirling coefficients in every order. & \pageref{sec:stirling-edge}\\[4pt]
8 & Booker--Str\"ombergsson--Then, Conj.~1 (2013) & True in every parameter-derivative order. A two-weight extension gives a necessary and sufficient parameter range. & \pageref{sec:jacobi}\\[4pt]
9 & Baricz--Wu, concluding remark~6 (2009) & False. Each proposed global extension fails. We determine the exact order range for the first two inequalities. & \pageref{sec:red}\\[4pt]
10 & Qi, Remark~7.4 (preprint 2013, journal 2020) & False. The exact replacement degrees are $2$, $3$, or a uniquely specified value strictly between them. & \pageref{sec:degrees}\\
\bottomrule
\end{tabular}\end{center}
\clearpage
\noindent\textbf{From auxiliary problems to the original questions.}
Three transfers are especially direct.
\begin{enumerate}
\item For partitions, ask how many input partitions can produce the same
pair of lists after a cut. A fixed finite subset of the allowed parts
gives many ways to fill both lists back to prescribed sums. A second
cut deals with the inputs that leave too little room on the first pass.
This gives a bound uniform in the larger input, including for sets
having arbitrarily large gaps.
\item For graph powers, ask whether adding a bounded part to a
composition of fixed total makes the sum of two distinguished parts
stochastically smaller. The comparison translates to the probability
that deleting a random member of a dominating set preserves domination.
That probability gives the ratio of successive normalized coefficients.
\item For the Kanter factor, ask when the argument of a boundary
value of its transform is increasing. The known probabilistic
criterion turns this into a two-parameter sign problem. A contour with
a uniformly controlled oscillatory integral deals with large arguments.
On the compact remainder, derivative bounds turn finitely many rounded
calculations into estimates on the entire parameter continuum.
\end{enumerate}

\noindent\textbf{Connection with the source corpus.}
We consulted the
\href{https://github.com/QEEDDD/Verifiable-Mathematical-Reasoning-Corpus/commit/ecd87de3348dcb706a76bb12fde722d7a348a012}{Verifiable Mathematical Reasoning Corpus, revision \texttt{ecd87de3348dc}}.
Item [076] uses a first crossing of a prefix deficit to reconstruct
combinatorial objects. Here the crossing of a sorted sum is used for
an estimate on the number of preimages. Items [067] and [083] provided
two other points of departure: converting distributional questions to
signs of transforms, and separating compact inequalities from bounds
at the boundary of an unbounded domain. The statements about the
Kanter factor and the curvature of the beta manifold are different.
The theorems in this manuscript use the mathematical inputs specified
in their respective proofs. Earlier results of the corpus have not been
counted as resolutions in the table.

\medskip
\noindent\textbf{Related graph problem and priority.}
The existence of ordinary domination polynomials with multiple peaks
had already been reported in 2026, including S. J. Silverman's July
\href{https://doi.org/10.5281/zenodo.21587251}{preprint}
and L. R. Chalmers's October
\href{https://doi.org/10.5281/zenodo.23092986}{formalization archive}.
The ordinary-domination conjecture is excluded from the table.
Entry~3 concerns powers of two specified graph families and a stronger
inequality. Entry~4 compares modes for pairs of unimodal trees.
Boyer et al.'s zero-forcing conjecture also has a prior counterexample:
P.~Graham, \emph{A graph whose zero-forcing polynomial is not
unimodal}, \href{https://doi.org/10.5281/zenodo.23875153}{1 November
2026}. Appendix~\ref{sec:zfjoin} gives a smaller graph and an
arbitrary-number-of-peaks construction. This is outside the ten.

\medskip
\noindent\textbf{Verification.}
The source contains complete listings for the square-partition finite
check, the independent zero-forcing enumeration, and the interval checks
for the Kanter and trigamma results. A fifth verifier uses rational
polynomials for the shifted-law obstruction. The finite-bell-shape results in the appendices distinguish exact
printed zero counts from the older endpoint convention. Under that
convention relevant positive implications were already proved by
Kwa\'snicki and Simon. They are excluded from the ten. Appendix~\ref{sec:lms-spectral} gives a further application of
a fractional-kernel comparison from the previous project.
\clearpage


\section{Submultiplicativity for partitions with arbitrary allowed parts}
\label{sec:arbitrary-A}
For $A\subseteq\mathbb{N}$ let $P_A(j)$ count unordered partitions of $j$
with all parts in $A$. Distinct partitions record only multiplicities.
Conjecture~6.5 in the 2024 preprint of K. Gajdzica, \emph{On the Bessenrodt--Ono type
inequality for a wide class of $A$-partition functions},
\href{https://arxiv.org/abs/2401.16267}{arXiv:2401.16267},
published in \href{https://doi.org/10.1007/s11139-025-01048-z}
{Ramanujan J. 66 (2025), article 80},
asks if $P_A(h)P_A(n)>P_A(h+n)$ for all sufficiently large $h,n$,
for every finite or infinite $A$ with $|A|\ge2$ and $\gcd A=1$.
There is no regularity assumption on the spacing of the parts.
\begin{theorem}\label{th:arbitrary-A}
Under exactly these hypotheses,
\[
 \lim_{H\to\infty}\ \sup_{n\ge h\ge H}
 \frac{P_A(h+n)}{P_A(h)P_A(n)}=0.
\]
Furthermore the convergence is uniform over all $A$ containing a
prescribed finite set $C$ with $\gcd C=1$ and $|C|\ge5$.
\end{theorem}

\begin{proof}
For a fixed finite set $C$ of $r\ge2$ parts and greatest common divisor
one, $p(j)=P_C(j)$ has the estimate
\[
 p(j)=\frac{j^{r-1}}{(r-1)!\prod_{a\in C}a}+O_C((1+j)^{r-2}).
\tag{A1}
\]
In fact $\prod_{a\in C}(1-z^a)^{-1}$ has its unique pole of
order $r$ at $1$, with leading coefficient $1/\prod a$ in powers of
$(1-z)^{-1}$. All other poles are roots of unity of order at most
$r-1$ as poles, so partial fractions give (A1).
For finite $A$ this already proves the result, since $n+h\le2n$.
For infinite $A$ take a finite gcd-one subset $C\subset A$ and augment
it until $r=|C|\ge5$. Such a subset exists by stabilization of the
gcd of the initial segments of $A$. We prove the stronger uniform
claim with this $C$ held fixed.

\paragraph{Counting the originals of a cut.}
Let $D=A\setminus C$. We will assign a partition of $n+h$ to a pair
$(L,R)$ of $D$-partitions, called its prefix and suffix, leaving
at most three individual $D$-parts unassigned. Define
$u=h-|L|$, $v=n-|R|$, where bars denote sums of parts. All pairs
of full partitions of $(h,n)$ with precisely these two $D$-parts
multisets are counted by $p(u)p(v)$. As $(L,R)$ varies these
classes are disjoint and their sizes sum to $P_A(h)P_A(n)$.

Suppose first that the original $C$-parts have total weight
$s\le2\min(u,v)$ and no more than three $D$-parts were removed.
For fixed $(L,R)$ and $s$ the missing parts sum to $u+v-s$.
There are $O((1+u+v)^2)$ possibilities, by counting ordered
compositions into at most three positive terms and allowing an
empty composition. The original $D$-multiset is the union of the
known multisets and the missing parts. There are $p(s)$ choices of
its $C$-multiset. Write $x=\min(u,v)$, $y=\max(u,v)$. By (A1) the
number of possible originals in this class is at most
$K_C(1+x)^r(1+y)^2$. For $x,y$ sufficiently large the ratio to
the class size $p(u)p(v)$ is therefore at most
\[
 K'_C x\,y^{3-r}\le K'_C x^{4-r}.
\tag{A2}
\]
In a second type of assignment no part is deleted and all $D$-parts
go on the right, with $s\ge2h$. Here $u=h$, $v=s-h\ge h$ and
the number of possible originals is bounded by $p(s)=p(h+v)$.
For sufficiently large $h$ the ratio to $p(h)p(v)$ is at most
$K''_C h^{1-r}$, again by (A1). Both estimates allow arbitrary $D$.

\paragraph{A cut preserving space for the finite set.}
Fix $\varepsilon>0$ and choose an integer $B$ so large that (A2)
is at most $\varepsilon/4$ whenever $x\ge B$.
Take $h\ge 4B$ sufficiently large that the second type of assignment
also costs at most $\varepsilon/4$ of a class. Each partition with
$s\ge2h$ receives precisely this second assignment.

For $s<2h$ order all $D$-parts nondecreasingly. Their total
$h+n-s$ exceeds $h-s/2>0$. Remove the first part $q$ for which
the cumulative sum reaches $h-s/2$. If $L_0,R_0$ are the preceding
and following lists, respectively, set $u_0=h-|L_0|$,
$v_0=n-|R_0|$. The definition of the crossing gives
\[
    s/2<u_0\le s/2+q,\qquad v_0=s+q-u_0\ge s/2.
\tag{A3}
\]
If $u_0<B$, remove the largest remaining part of $L_0$. If $v_0<B$,
remove the smallest remaining part of $R_0$. Each list from which
a removal is required is nonempty: an empty $L_0$ would give
$u_0=h\ge B$, and an empty $R_0$ would give $v_0=n\ge B$.
After these operations at most three parts have been removed.
If the new capacities $u,v$ are both at least $B$, use the new
prefix and suffix as $(L,R)$. Inequality (A3) ensures
$s\le2\min(u,v)$, so the first estimate applies.
Thus the total number of partitions assigned so far is at most
$(\varepsilon/2)P_A(h)P_A(n)$, by summing the two estimates over
all classes.

\paragraph{The partitions left over.}
We show that every unassigned partition has at least $h-2B$ of its
weight in $C'=C\cup(A\cap[1,B])$. If the final left capacity is
below $B$, the largest part of $L_0$ was removed and was less than
$B$. Thus the entire list $L_0$ consists of parts less than $B$,
of total $h-u_0>h-B$. If instead the final right capacity is below
$B$, the smallest part of $R_0$ was less than $B$, so $q<B$ by
the ordering. Also $s/2\le v_0<B$. Now $u_0\le s/2+q<2B$
by (A3), and every part of $L_0$ is at most $q<B$. In this case
$|L_0|>h-2B$, as required.

Apply the construction again to these original partitions, with their
original parts intact and the new finite set $C'$ in place of $C$.
Write $s'$ for its weight in a partition. If $s'\ge2h$, use the
assignment with every part outside $C'$ on the right. If $s'<2h$,
make the sorted cut and remove only its crossing part.
By (A3) both capacities are at least $s'/2\ge(h-2B)/2$.
Consequently the total number of partitions assigned in this second
pass is at most $(\varepsilon/2)P_A(h)P_A(n)$ for all sufficiently
large $h$, by the two preceding estimates for $C'$. There are
only finitely many possibilities for $C'$ once $C,B$ are fixed,
so the threshold can be taken independent of the ambient set $A$.
Adding the two counts proves the theorem. In particular the strict
inequality in Conjecture~6.5 follows by taking $\varepsilon<1$.
\end{proof}

Gajdzica's original finite reduction (Theorem~1.4) imposes separation
on the allowed parts and verifies several special sets. The sorted split
of an individual partition has precedents in combinatorial proofs of
Bessenrodt--Ono inequalities, including that source. Here the number
of preimages is compared with the number of ways to fill both sums
from a finite set. The second cut deals with every failure of the
first by using only bounded part sizes. This supplies the arbitrary-set
statement and the uniform vanishing ratio.

\begin{corollary}\label{th:arbitrary-colored}
Suppose the weight $j$ has $c(j)$ distinguishable colors, with
$0\le c(j)\le M(1+j)^\beta$ for fixed $M,\beta\ge0$.
Allow each colored part arbitrary multiplicity. If the occupied
weights have gcd one and there are at least two colored part types,
the same vanishing ratio holds for the colored partition counts.
\end{corollary}
\begin{proof}
For finitely many colored part types the finite-product proof of
(A1) applies with multiplicities. Otherwise choose a gcd-one finite
set of $r>4+3\beta$ colored part types. Order all remaining types by
weight with a deterministic order for ties. Each deleted part has
at most $M(1+u+v)^\beta$ colors. With at most three deletions, the
upper bound preceding (A2) gains a factor $O_M((1+y)^{3\beta})$.
Its ratio to class size is $O_{C,M}(x y^{3+3\beta-r})
=O_{C,M}(x^{4+3\beta-r})$, tending to zero. There are finitely
many colored types up to $B$, so the second cut and the rest of the
argument apply unchanged.
\end{proof}

\begin{corollary}\label{th:arbitrary-rate}
For every finite gcd-one set $C$ of $r\ge4$ distinct part sizes,
uniformly in $A\supseteq C$ and $n\ge h\to\infty$,
\[
 \frac{P_A(h+n)}{P_A(h)P_A(n)}=O_C\bigl(h^{-(r-3)/2}\bigr).
\]
Consequently, for each infinite $A$ of gcd one and every $K>0$
the quotient is $O_{A,K}(h^{-K})$ uniformly in $n\ge h$. For
a finite gcd-one set $A$ of $r\ge2$ part sizes the corresponding
supremum over $n\ge h$ is of order $h^{1-r}$.
\end{corollary}
\begin{proof}
In the first cut use $B=\lfloor\sqrt h\rfloor$ and put aside every
case where both $u_0,v_0<B$. In such a case
$q=u_0+v_0-s<2B$ and $|L_0|>h-B$. The successful cuts now have
at most two deleted parts. The upper count preceding (A2) drops to
$O_C((1+x)^r(1+y))$ because two missing weights of prescribed sum
have only $O(1+y)$ possibilities. Dividing by $p(u)p(v)$ gives a
cost $O_C(xy^{2-r})=O_C(B^{3-r})$. The first-pass no-deletion
case still costs $O_C(h^{1-r})$.

For every partition put aside, the parts in $L_0$ have weight greater
than $h-2B$ and each has size at most $2B$. This was proved above
for an unsuccessful single-side removal, and it holds by construction
when both sides are small. One size $d\le2B$ in $D$ must therefore
carry at least $(h-2B)/(2B)\ge h/(4B)$ of the original weight for
large $h$. Choose one such size by a fixed tie-breaking rule.
Use $C_d=C\cup\{d\}$ alone for this partition's second cut.
The counts $p_d(t)=P_{C_d}(t)$ satisfy
\[
 p_d(t)\asymp_C t^{r-1}(1+t/d)\qquad(t\ge t_0(C)),
\]
with constants independent of $d$. Sum $p(t-ld)$ for
$0\le l\le t/d$, using $0\le l\le t/(2d)$ for the lower bound.
The analogous upper bound with $1+t$ holds for all $t\ge0$.

If the new core weight is at least $2h$, the no-deletion estimate
remains $O_C(h^{1-r})$ for each $d$, since
$(1+(h+v)/d)\le(1+h/d)(1+v/d)$.
Otherwise the second cut deletes only its crossing, so for fixed
$s,L,R$ the deleted weight is uniquely determined. Both remaining
capacities are at least $h/(8B)$. The class count is now at most
$O_C((1+x)^r(1+x/d))$, by summing $p_d(s)$ for $s\le 2x$.
Division by $p_d(u)p_d(v)$ costs $O_C(xy^{1-r})=O_C(x^{2-r})$.
There are at most $2B$ choices of $d$, giving the combined cost
\[
 O_C\bigl(B^{3-r}+Bh^{1-r}+B(h/B)^{2-r}\bigr)
 =O_C\bigl(h^{-(r-3)/2}\bigr).
\]
For infinite $A$ choose $r$ arbitrarily large. For finite $A$,
(A1) gives the uniform upper bound, and $n=h$ gives the matching
lower bound.
\end{proof}

\paragraph{A sharp boundary for universal color bounds.}
Let $W:\mathbb{N}\to[1,\infty)$ be a prescribed upper bound on the
number of colors at each weight. Every coloring $c(j)\le W(j)$
with gcd-one occupied weights and at least two colored part types
obeys eventual strict submultiplicativity if and only if $W$ is
bounded by a polynomial. Sufficiency was proved above. For necessity
start with one color at each of the weights $1,2$. Having chosen
finitely many additional colored weights, let $R$ be the current
total number of types. Their partition function $p_0(t)$ is
$O((1+t)^{R-1})$ by (A1). If $W$ has no polynomial upper bound,
there are arbitrarily large $N$ at which
\[
 \lfloor W(N)\rfloor>
 j\,p_0(\lfloor N/2\rfloor)p_0(\lceil N/2\rceil)
\]
at step $j$. Choose such an $N$ greater than twice the largest
previously selected weight, and introduce $\lfloor W(N)\rfloor$
colors at $N$. For the resulting infinite colored partition function
the denominator at the two halves is unchanged forever, whereas
the numerator at $N$ counts at least all the new single-part
partitions. The ratio at these balanced pairs therefore tends to
infinity along the selected weights. No extra named conjecture is
counted for this consequence.

\begin{proposition}\label{th:three-core}
Three distinct prescribed parts of gcd one suffice for convergence
to zero uniformly over all sets containing them, with bound
$O_C(h^{-1/5})$ for $n\ge h$. For every prescribed two-element
gcd-one set, eventual strict submultiplicativity fails to have
a threshold uniform over such supersets.
\end{proposition}
\begin{proof}
For a three-part core modify the quantitative proof with
$B=\lfloor h^{1/5}\rfloor$, $J=\lfloor h^{2/5}\rfloor$.
A successful cut with only its crossing deleted costs
$O_C(xy^{1-r})=O_C(B^{-1})$ at $r=3$. A two-deletion cut must
have had one initial capacity below $B$, whence $s<2B$.
If its final $y=\max(u,v)>J$, its preimage count is
$O_C(B^3 y)$, compared with $p(u)p(v)\gg_C x^2y^2$; the
relative cost is $O_C(B/J)$. If instead $y\le J$, we have
$q\le u_0+v_0\le2J$. All of $L_0$ has parts at most $2J$ and
weight $h-u_0\ge h-J$. Put this case aside.
The other unsuccessful cuts still have at least $h-2B$ of weight
on parts in $L_0$ of sizes at most $2B$. For all partitions put
aside, some size $d\le 2J$ outside $C$ therefore carries at least
$h/(4J)$ of weight for large $h$. The single-size second pass
then costs $O_C(J(h/J)^{-1}+Jh^{-2})$. The total cost is
$O_C(B^{-1}+B/J+J^2/h)=O_C(h^{-1/5})$.

For the opposite direction fix $C=\{a,b\}$, $a<b$, with
$\gcd(a,b)=1$, and for each large $h$ set
$A_h=C\cup\{h+1,\ldots,2h\}$. With $p=P_C$,
\[
 P_{A_h}(h)=p(h),\qquad
 P_{A_h}(2h)=p(2h)+\sum_{j=0}^{h-1}p(j).
\]
Since $p(j)=j/(ab)+O_C(1)$, the quotient at $(h,h)$ tends
to $ab/2$, exceeding one unless $C=\{1,2\}$. For this last set
use $p(j)=\lfloor j/2\rfloor+1$. At $h=2m$ the difference
$P_{A_h}(2h)-P_{A_h}(h)^2$ is exactly $m>0$.
\end{proof}

\section{The exception table for multiplicativity of square partitions}
\label{sec:roy-squares}
Write $P_A(n)$ for the number of unordered partitions of $n$ with parts
in a set $A$ of positive integers, with $P_A(0)=1$ and $P_A(n)=0$ for
$n<0$. For $A=\{1,4,9,\ldots\}$ write $P(n)=P_A(n)$. In
\emph{Log-concavity and the multiplicative properties of restricted
partition functions}, \href{https://doi.org/10.1007/s00026-026-00847-5}
{Ann. Comb. (2026)}, Conjecture~3.2, Arindam Roy gives an exact list of
the equalities and failures of $P(h)P(n)\ge P(n+h)$ for $n\ge h\ge1$.
The preprint dates to 2024 (\href{https://arxiv.org/abs/2404.03153}
{arXiv:2404.03153}). Gajdzica, \emph{On the Bessenrodt--Ono type
inequality for a wide class of $A$-partition functions}, Theorem~5.2
(\href{https://arxiv.org/abs/2401.16267}{arXiv:2401.16267}), had already
proved strict inequality for $h,n\ge12$. This leaves infinitely many
comparisons with one argument below $12$.

\paragraph{An auxiliary extension problem.}
Let $C$ be a finite set of allowed parts and let $h$ be fixed.
What conditions preserve a shift inequality if arbitrarily many large
parts are added? For integer $s\ge-h$, put
$g(s)=P_C(h)P_C(s)-P_C(s+h)$. Suppose the negative mass
$B=\sum_{s\ge-h}\max\{-g(s),0\}$ is finite. If an integer $J>h$ satisfies
$g(s)>B$ for all $s\ge J-h$, then for every
$D\subseteq\{J,J+1,\ldots\}$ disjoint from $C$ and all $n\ge J-h$,
\[
 P_{C\cup D}(h)P_{C\cup D}(n)>P_{C\cup D}(n+h).
\tag{R}
\]
Indeed $P_{C\cup D}(h)=P_C(h)$ and the difference in (R) is the sum
of $g(n-|\lambda|)$ over all $D$-partitions $\lambda$ of weights up
to $n+h$. A negative term must have nonempty $\lambda$, since the
empty partition gives $g(n)>B$. Delete its smallest part $d$. The
resulting $\lambda'$ contributes $g(s+d)$ where $s=n-|\lambda|\ge-h$.
As $s+d\ge J-h$, this contribution exceeds $B$. For fixed $\lambda'$
and fixed $s$, the deleted part is forced to be $n-|\lambda'|-s$.
Thus the sum of all deficits assigned to one positive term is at most
$B$. This leaves the empty-partition term strictly positive after all
assignments and proves (R).

\begin{theorem}\label{th:roy-square}
Roy's Conjecture~3.2 holds. In fact, in the eight strips $4\le h\le11$,
the complete exception table remains valid when all squares at least $49$
are replaced by any set of positive integer parts at least $49$.
\end{theorem}
\begin{proof}
Apply (R) to $C=\{1,4,9,16,25,36\}$ and $J=49$. For $n+h<49$ the
counts are those of $C$ itself. The following table gives the negative
mass $B$, the minimum of $g$ on $s\ge49-h$, and the complete non-strict
cases for $n\ge h$ in each strip.
\[
\begin{array}{c|rr|l|l}
 h&B&\min g& P(h)P(n)=P(n+h)&P(h)P(n)<P(n+h)\\ \hline
 4&4&59&5,6,7&\varnothing\\
 5&8&45&5,6,8,11,15&7\\
 6&16&35&8,10,12,14&6,7,11,15\\
 7&26&29&8,9,12,13,19&7,10,11,14,15\\
 8&15&83&\varnothing&\varnothing\\
 9&16&123&\varnothing&\varnothing\\
 10&23&103&\varnothing&\varnothing\\
 11&31&97&\varnothing&\varnothing
\end{array}
\]
Here is a finite certificate for the infinite statements about $g$.
Let $L=3600$ and compute $p[0],\ldots,p[6L+20]$ by starting with
$p[0]=1$ and zeros elsewhere. For each $a\in C$, in increasing
order of $n$ from $a$ to $6L+20$, add $p[n-a]$ to $p[n]$.
This gives $p[n]=P_C(n)$ by multiplication of generating series.
For each residue $0\le r<L$, $Q_r(k)=g(r+kL)$ is a polynomial in $k$
of degree at most five on all $k\ge0$. Indeed,
\[
 \prod_{a\in C}(1-z^a)^{-1}
 =(1-z^L)^{-6}\prod_{a\in C}(1+z^a+\cdots+z^{L-a}).
\]
The numerator has degree less than $6L$. The contribution of a
monomial in residue class $r$ involves $\binom{k-j+5}{5}$ with
$0\le j\le5$. When $0\le k<j$ this polynomial is zero, so validity
extends all the way to $k=0$. The shift by $h$ in $g$ preserves this
polynomiality. Compute $Q_r(0),\ldots,Q_r(5)$ from the $p[n]$.
Successively replacing adjacent entries by their forward differences
gives $\Delta^jQ_r(0)\ge0$ for $1\le j\le5$ for every $h,r$ listed.
The binomial-basis expansion of $Q_r$ therefore makes it nondecreasing
for all $k\ge0$. Its values on $L\le s<2L$ all exceed the displayed
$B$. It now suffices to sum the deficits for $-h\le s<L$ and take
the minimum for $49-h\le s<2L$ to reproduce the two columns.
All these exact integer operations are provided in the accompanying
short verifier, \texttt{roy\_square\_verify.py}.

For $h=1,2,3$ one has $P(h)=1$. The increment $P(j)-P(j-1)$
counts the square partitions of $j$ lacking a part $1$. Every
$j\ge24$ is a sum of fours and nines: the representatives
$0,9,18,27$ for the four residue classes suffice. Consequently
$P(n+h)>P(n)$ for all $n\ge23$. A direct recurrence up through $25$
gives equality, for $h=1$, at $n=1,2,4,5,6,9,10,13,14,18,22$,
for $h=2$ at $n=4,5,9,13$, and for $h=3$ at $n=4$.
Every other pair in these strips fails. Gajdzica's theorem covers
$h\ge12$. Comparing these disjoint ranges with the published table
proves the conjecture.
\end{proof}

\section{Ultra-log-concavity for dominating sets in powers of paths and cycles}
\label{sec:gz-ulc}
Galvin and Zhang, \emph{The domination polynomial of powers of paths
and cycles}, \href{https://arxiv.org/abs/2408.12731}{arXiv:2408.12731},
prove unimodality and ask in Conjecture~1.8 whether both coefficient
sequences are ultra-log-concave for every order $n$ and power $\ell\ge1$.
For a graph on $n$ vertices set $f_k=d_k/\binom nk$, where $d_k$ counts
its $k$-element dominating sets. Their conclusion asks for
$f_k^2\ge f_{k-1}f_{k+1}$, not merely log-concavity of $d_k$.

\begin{theorem}\label{th:gz-ulc}
For every $\ell\ge1$, the ordinary domination sequences of $P_n^\ell$
and $C_n^\ell$ are ultra-log-concave wherever the corresponding
graphs are defined.
\end{theorem}

\paragraph{A comparison between bounded compositions.}
Let $U_b(x)=x+x^2+\cdots+x^b$ for $b\ge1$.
Fix an integer $N$ and a nonnegative finite sequence $v=(v_j)$.
Consider a random integer with probabilities proportional to
$v_j h_{N-j}$, where $h$ is the coefficient sequence of a product
of polynomials $U_b$, allowing the empty product. Replacing $h$ by
the coefficients $h'= h*[x^\bullet]U_c$ for any $c\ge1$ makes
the random integer stochastically smaller, whenever both distributions
can be normalized.

Indeed $h$ is log-concave with interval support. To recall the
convolution argument, the Toeplitz matrix $(h_{u-v})_{u,v\in\mathbb Z}$
has all order-two minors nonnegative exactly when its interval-supported
sequence is log-concave (multiply consecutive ratio inequalities).
Interval indicators have this property. Their convolutions do too
by matrix multiplication and the order-two Cauchy--Binet identity;
finite support makes the sums finite. The delta mass at zero also has
the property. For $u_1<u_2$ we consequently have
\[
 h'_{u_2}h_{u_1}-h'_{u_1}h_{u_2}
 =\sum_{i=1}^{c}\{h_{u_2-i}h_{u_1}-h_{u_1-i}h_{u_2}\}\ge0.
\]
This comparison includes the zeros beyond the supports. Put $u=N-j$.
For $j\le t<j'$ it gives
$(v_jh'_{N-j})(v_{j'}h_{N-j'})\ge
(v_jh_{N-j})(v_{j'}h'_{N-j'})$. Summing over these pairs and
normalizing shows that the probability of $j\le t$ increases.

\paragraph{Deletion probabilities.}
In any nonempty graph the indices with $f_k>0$ form a terminal interval.
Given a uniform dominating $k$-set choose one of its members uniformly
and delete it. The probability that the remaining set dominates is
\[
 s_k=\frac{(n-k+1)d_{k-1}}{k d_k}=\frac{f_{k-1}}{f_k}.
\]
Every dominating $(k-1)$-set indeed has all $n-k+1$ possible
one-element extensions. It suffices to show $s_k$ increases weakly
throughout the positive range. For $k=1$ the deletion probability
is zero. Empty graphs have a single coefficient.

Put $Q=2\ell+1$. On the cycle choose a distinguished selected vertex
uniformly and list the successive clockwise distances between selected
vertices, starting there. The ordered $k$ distances are uniform among
compositions of $n$ with all parts in $[1,Q]$: each feasible ordered
tuple and starting vertex gives exactly one pair of a dominating set
and its distinguished member. Bounding each part by $Q$ is equivalent
to covering every vertex between consecutive selected vertices. For
$k\ge2$, deleting the distinguished member preserves domination exactly
when the sum $j$ of the two incident gaps is at most $Q$. Its law has
weights
\[
 ([x^j]U_Q^2)([x^{n-j}]U_Q^{k-2}).
\]
Going from $k$ to $k+1$ multiplies the residual polynomial by $U_Q$.
The comparison just proved makes $j$ stochastically smaller. Hence
$s_{k+1}\ge s_k$ on the cycle, including the case $k=2$.

On the path with vertices $1,\ldots,n$, for selected positions
$x_1<\cdots<x_k$ form the end gaps $x_1,n+1-x_k$ and interior gaps
$x_{i+1}-x_i$. Their sum is $n+1$. Domination is equivalent to both
end gaps being at most $A=\ell+1$ and all interior gaps being at most
$Q$. These positive gap vectors are uniform at fixed $k$ since each
gives one subset. For $k\ge2$ let $u_k$ be the probability that the
first selected vertex can be deleted. This is the probability that
the first end gap plus the first interior gap is at most $A$. The
last selected vertex has the same probability by reflection. For
$k\ge3$ let $v_k$ be the corresponding probability for an interior
selected vertex. Exchangeability of the interior gaps makes this
independent of the interior vertex index. Its condition is that two
adjacent interior gaps sum to at most $Q$.

For each type of pair, its weight at a sum $j$ is the coefficient at
$j$ of the product of its two gap polynomials, multiplied by the
coefficient at $n+1-j$ of all the other gap polynomials. Passing from
$k$ to $k+1$ multiplies only the latter product by $U_Q$.
Consequently $u_{k+1}\ge u_k$ for $k\ge2$, and $v_{k+1}\ge v_k$ for
$k\ge3$ whenever the probabilities are defined.

There is also a comparison at the same $k\ge3$. Write $a$ for the
left end gap, and $b,c$ for the first two interior gaps. Conditional
on $b$, all the other gaps apart from $a,c$, and the resulting sum
$a+c=S$, the two-part split is uniform over $1\le a\le A$,
$1\le c\le Q$. The random variables $A-a$ and $Q-c$ are uniform on
integer intervals of equal length, since $a\mapsto c=S-a$ is a
bijection. Their least possible values are
$\max(0,A+1-S)$ and $\max(0,Q+1-S)$, respectively. Since $Q\ge A$,
$Q-c$ stochastically dominates $A-a$. Comparing the probabilities
that their values are at least $b$ proves
\[
 \mathbb P(b+c\le Q)\ge\mathbb P(a+b\le A),\qquad v_k\ge u_k .
\]
At $k=2$ one has $s_2=u_2$. For $k\ge3$ one has
$s_k=2u_k/k+(1-2/k)v_k$. Both probabilities increase with $k$ and
the weight of the larger one increases, so $s_{k+1}\ge s_k$.
At the transition to $k=3$ the same conclusion follows from
$s_3\ge u_3\ge u_2$. This proves Theorem~\ref{th:gz-ulc}.

\section{Deleting one leaf can move a tree's domination mode arbitrarily far}
\label{sec:tree-leaf-mode}
S. Zhang, \emph{Domination polynomial is unimodal for large graphs
with a universal vertex}, \href{https://arxiv.org/abs/2111.00641}{arXiv:2111.00641},
Conjecture~4.1, predicts that deleting a leaf changes the mode of a
tree's domination polynomial by at most one whenever both polynomials
are unimodal. This is the all-tree question of Burcroff and O'Brien,
\href{https://arxiv.org/abs/2110.00709}{arXiv:2110.00709}, Question~8.1.
Their separate Conjecture~3.7 in that preprint concerns short-legged
spiders.

\begin{theorem}\label{th:tree-leaf-mode}
There are trees $T$ and leaves $v$ for which both ordinary domination
polynomials are unimodal and every mode of $D(T-v,x)$ exceeds every
mode of $D(T,x)$ by an arbitrarily large amount.
\end{theorem}

\paragraph{The tree and its exact polynomial.}
For $d\ge1$ put $K=2\cdot3^d$. At a vertex $r$ attach $b$ private
leaves, where $b=1,2$, and $K$ neighbors $u_i$. Attach to each $u_i$
$d$ disjoint paths $u_i-t_{ij}-w_{ij}$ of length two. Write $T_b$
for this tree. Deleting a private leaf of $r$ from $T_2$ gives $T_1$.
With
\[
 A(x)=2x+x^2,\qquad P(x)=(1+x)A(x)^d,\qquad Q(x)=P(x)-x^d,
\]
the polynomial is exactly
\[
 D_b(x)=x(1+x)^bP(x)^K+x^b Q(x)^K.
\]
Indeed at least one vertex of each pair $\{t_{ij},w_{ij}\}$ must be
selected. If $r$ is selected the choice of $u_i$ is free. Otherwise
all $b$ private leaves are selected and dominate $r$; the excluded
choice in each branch is precisely to omit $u_i$ and all $t_{ij}$
and select all $w_{ij}$.

For any nonzero polynomial $L$ with nonnegative coefficients, let
$X_L$ denote the integer random variable with probability generating
function $L(x)/L(1)$. The law of $X_P$ is
$d+\mathrm{Bin}(d,1/3)+\mathrm{Bernoulli}(1/2)$ with independent
terms. The law of $X_Q$ results from removing mass $1/K$ at $d$ and
renormalizing. Thus for $m_P=\mathbb E X_P$ and $m_Q=\mathbb E X_Q$,
\[
 m_P=\tfrac43d+\tfrac12,\qquad
 m_Q-m_P=\frac{d/3+1/2}{K-1}.
\]
The two summands of $D_b(1)$ have masses proportional to $2^b$ and
$h_K=(1-1/K)^K$ respectively. Writing $\alpha_b=h_K/(2^b+h_K)$
gives the mean size $\mu_b$ under the uniform law on dominating sets:
\[
 \mu_b=(1-\alpha_b)(Km_P+1+b/2)+\alpha_b(Km_Q+b).
\]
In particular
\[
 \mu_1-\mu_2=
 \frac d3\left(\frac{e^{-1}}{2+e^{-1}}-
               \frac{e^{-1}}{4+e^{-1}}\right)+O(1).
 \tag{*}\label{leaf:gap}
\]
Below we show that $D_b$ is unimodal and every mode is within $O(1)$
of $\mu_b$, uniformly as $d\to\infty$.

\paragraph{The two summand distributions.}
Let $p(k)$ stand for the normalized coefficients of either summand.
Its law is a sum of $K$ independent branch sizes ($X_P$ or $X_Q$)
and one independent bounded variable encoding the root and private leaves.
Write $\nu,V,\kappa_3$ for its mean, variance and third cumulant.
For each fixed $C>0$, uniformly when $|k-\nu|\le Cd$,
\begin{align}
 \sqrt{2\pi}V^{3/2}\,[p(k+1)-p(k)]
 &=\nu-k-\tfrac12-\frac{\kappa_3}{2V}+o(1),\label{leaf:first}\\
 \sqrt{2\pi}V^{3/2}\,[p(k+1)-2p(k)+p(k-1)]&=-1+o(1).
 \label{leaf:second}
\end{align}
In both distributions
\[
 V=V_0+O(d^2),\qquad V_0=K(2d/9+1/4),\qquad |\kappa_3|=O(Kd).
\]
To see this, removing probability $1/K$ from a law supported on
$[d,2d+1]$ changes its first three raw moments by $O(d^j/K)$
at order $j$. Hence its second and third cumulants change by
$O(d^2/K)$ and $O(d^3/K)$, respectively. The Bernoulli sum for $X_P$
has variance $2d/9+1/4$ and third cumulant $O(d)$; cumulants add.

We give details of the local estimates because a weak normal limit
alone does not locate a mode to bounded precision. The single-branch
characteristic function for $X_P$ has modulus at most
$e^{-c d t^2}$ on $|t|\le\pi$, by
$|1-p+pe^{it}|^2=1-2p(1-p)(1-\cos t)$.
For $X_Q$ the modulus is at most $(K e^{-c d t^2}+1)/(K-1)$.
Its $K$th power is bounded by
$C_1e^{-c_1K\min(dt^2,1)}$: factor out $(K+1)/(K-1)$
and use $1-e^{-cdt^2}\ge c_2\min(dt^2,1)$. The bounded
independent factor has characteristic modulus at most one. Thus the
Fourier integral over $K^{-2/5}<|t|\le\pi$ is negligible even
after multiplication by any fixed power of $K$.

On the remaining interval, the characteristic function of the
centered sum has expansion
\[
 e^{-Vt^2/2}\left(1+\frac{\kappa_3(it)^3}{6}
       +O\big(Kd^4 |t|^4+(Kd)^2|t|^6\big)\right).
\]
For justification, a centered branch is bounded by $2d+1$.
On a complex disk of radius $c_3/d$ about zero its characteristic
function is within $1/2$ of one for sufficiently small $c_3$.
The analytic logarithm is bounded there. Its Taylor remainder after
order three is $O(d^4|t|^4)$ by Cauchy's estimate on half the disk.
Sum these logarithms, including the bounded independent factor,
and exponentiate. On $|t|\le K^{-2/5}$ the cubic term and the
fourth-order remainder both tend uniformly to zero.

Inverting for the forward difference introduces the multiplier
\[
 (e^{-it}-1)e^{-it(k-\nu)}
 =-it-(k-\nu+1/2)t^2+O\big((1+|k-\nu|^2)|t|^3\big).
\]
The odd leading term integrates to zero against the centered
Gaussian. Its product with the cubic term contributes
$-\kappa_3t^4/6$. Integrating the remaining even terms and using
$\int_{\mathbb R}t^4 e^{-Vt^2/2}\,dt/
\int_{\mathbb R}t^2 e^{-Vt^2/2}\,dt=3/V$ yields
\eqref{leaf:first}. For example, after normalization the multiplier
remainder is $O((1+d^2)/\sqrt V)$ and the characteristic remainder
times $O(|t|)$ contributes
$O(Kd^4/V^{3/2}+(Kd)^2/V^{5/2})$, both $o(1)$.
The product of the quadratic multiplier term and the cubic
characteristic term is likewise $O((1+d)Kd/V^{3/2})=o(1)$
after normalization. The omitted Gaussian tails are exponentially small.
The central second difference uses
$-t^2+O((1+|k-\nu|)|t|^3)$ as multiplier; the same bounds give
\eqref{leaf:second}.

\paragraph{Unimodality and the mode displacement.}
Each summand has log-concave coefficients and interval support.
Indeed $P$ is, up to a monomial, a product of positive linear
polynomials, and $Q$ differs only in a reduced lowest nonzero
coefficient; this reduction preserves log-concavity and positivity
of that coefficient. Convolution preserves log-concavity of
nonnegative interval-supported sequences, as shown in the proof of
Theorem~\ref{th:gz-ulc} (or by the order-two Toeplitz minors).
Consequently \eqref{leaf:first} brackets all modes of each summand
between its mean minus and plus an absolute constant, by checking
the sign of the first difference on either side.

The two means differ by $O(d)$. Below a common central interval
containing both mode ranges the mixture increases, and above it
the mixture decreases, since its two summands do so individually.
Throughout that interval the mixture is strictly concave in the
discrete sense by \eqref{leaf:second}, proving unimodality.
To locate its own modes, multiply its forward difference by
$\sqrt{2\pi}V_0^{3/2}$. Applying \eqref{leaf:first} to the two
summands and averaging gives $\mu_b-k+O(1)$ uniformly there.
Indeed $V_0^{3/2}/V^{3/2}=1+O(d^2/V_0)$, and its error even when
multiplied by $O(d)$ tends to zero. Thus every mode is $\mu_b+O(1)$.
Equation~\eqref{leaf:gap} now proves the theorem.

\section{The generalized gamma convolution range of the Kanter variable}
\label{sec:kanter}
For $0<\beta<1$, the Kanter variable $K_\beta$ is the positive
factor in the independent factorization $Z_\beta\stackrel{d}{=}
L^{1-1/\beta}K_\beta$, where $L$ is unit exponential and
$\mathbb Ee^{-sZ_\beta}=e^{-s^\beta}$. Equivalently its Mellin
transform is given in (5) of Hasebe, Simon and Wang, \emph{Some
properties of the free stable distributions}, Ann. Inst. H. Poincar\'e
Probab. Statist. 56 (2020), 296--325,
\href{https://arxiv.org/abs/1805.01133}{arXiv:1805.01133}.
A generalized gamma convolution (GGC) is a weak limit of finite sums
of independent gamma variables with arbitrary scales.
\begin{theorem}\label{th:kanter}
The law of $K_\beta$ is a GGC if and only if $\beta\le4/5$.
Thus Conjecture~1 of Hasebe, Simon and Wang holds.
\end{theorem}
They prove necessity in Corollary~2 and sufficiency for $\beta\le3/4$
via Proposition~1. Their Thorin phase criterion (13)--(14) reduces
the outstanding interval to the following problem. Put $a=1-\beta$,
$1/5\le a\le1/4$ and $x=r^{1-a}>0$. Define
\[
 R(x)+iI(x)=\sum_{n\ge0}
    \frac{e^{i\pi n a}x^n}{n!\,\Gamma(1-na)}.
\]
The Thorin phase is $\pi^{-1}\arg(R+iI)$. Its nondecrease on the whole
positive half-line is equivalent, by the source, to the GGC property.
We prove $R>0$ and $(I/R)'>0$ throughout this strip.
The reflection formula gives, for positive $n$,
\[
 [x^n]R=\frac{\Gamma(na)\sin(2\pi na)}{2\pi n!},\qquad
 [x^n]I=\frac{\Gamma(na)\sin^2(\pi na)}{\pi n!}.
\]
In particular $I>0$. The positive coefficients will control $I$ in
the compact part of the argument. The oscillating coefficients of
$R$ require a different estimate for interpolation.

\paragraph{Integral and contour.}
Expanding exponentials under the gamma integral gives
\begin{align*}
 R&=1+\frac1{2\pi}\Im\int_0^\infty e^{-t}
        (e^{x e^{2\pi ia}t^a}-1)\,\frac{dt}{t},\\
 I&=\frac1{2\pi}\int_0^\infty e^{-t}
        (e^{xt^a}-\Re e^{x e^{2\pi ia}t^a})\,\frac{dt}{t}.
\end{align*}
In the first integral one can take the imaginary part of the
unsubtracted integrand with a lower cutoff tending to zero.
Use the principal power and deform to the path from $0$ to $iS$
and then horizontally to $+\infty+iS$. The small circular arc
for the unsubtracted integrand contributes $+1/4$ to $R$. The
vertical connector at the right endpoint tends to zero since
$t^a=o(|t|)$. On the vertical segment $t=iu$ the exponential
$e^{-t+x e^{2\pi ia}t^a}$ has modulus at most one, because
$\phi=5\pi a/2\in[\pi/2,5\pi/8]$. At $t=u+iS$ its real exponent
has derivative at most $-1+axS^{a-1}$. Thus when $axS^{a-1}<1/2$
the horizontal exponential has modulus at most $e^{-u/2}$.
Its contribution from the unsubtracted integrand, and that to the first
$x$ derivative, tend to zero as $S\to\infty$ (bounded by constants
times $S^{-1}$ and $S^{a-1}$). This argument justifies the full
vertical representation even at $a=1/5$ where it is oscillatory.

\paragraph{Unbounded range.}
Put $p=1/a$, $T=(ax)^{1/(1-a)}$, $v_0=T^a$. If $x\ge37$ then
$T=(x/p)^{p/(p-1)}\ge(37/5)^{5/4}>12$, since the expression
decreases in $p$ here. Let $q=\sin\phi>0.92$,
$A=Tq^{p/(p-1)}>10$, $v_*=v_0 q^{1/(p-1)}$, and
$D=-pA\cot\phi\ge0$. On the vertical ray substitute $t=iAz^p$.
The formula for $R$ becomes
\[
 R=\frac54+\frac p{2\pi}\Im\int_0^\infty
         e^{-Dz+iA(pz-z^p)}\,\frac{dz}{z}.
\]
Replacing $pz-z^p$ by $pz$ gives imaginary integral
$\arg(D+ipA)=\pi-\phi$, by the elementary Laplace sine integral
(including $D=0$ by a limit). The corresponding value of $R$ is $p/2$.
We give an explicit upper bound $0.8$ for the error in the imaginary
integral. Throughout, a one-sided oscillatory integral with monotone
phase derivative of magnitude at least $b$ has modulus at most
$2/(Ab)$ when the exponential is $e^{iA\,\mathrm{phase}}$.
Integration by parts proves this inequality. Multiplication by a
nonnegative decreasing amplitude costs only its maximum, by averaging
indicators of initial subintervals.

Take $c=3/5$ and $y=z-z^p/p$ on $0<z<c$. Compare the integrand there
with the linear-phase integral up to $y(c)$. The absolute discrepancy
of the measures is at most
\[
 -\log(1-c^{p-1}/p)
 +\frac{c^{p-1}}{e p(p-1)(1-c^{p-1}/p)}<0.065 .
\]
Indeed $0<y'/y\le1/z$ and the integral of $1/z-y'/y$ is the first
term. The damping error is bounded using
$|e^{-Dz}-e^{-Dy}|\le(z-y)/(ey)$ and $y'/y\le1/z$.
The remaining linear tail from $y(c)$ has modulus at most
$2/[pAy(c)]<0.09$.
For the nonlinear tail split at $7/8$ and $9/8$. The absolute bounds
for the three pieces are respectively
\[
 \frac{2}{c p A(1-(7/8)^{p-1})},\qquad
 \log(9/7),\qquad
 \frac{2}{(9/8)pA((9/8)^{p-1}-1)} .
\]
Their sum is below $0.62$, giving the claimed error. For $R'$ use the
differentiated contour before rescaling, which replaces $dz/z$ by
$e^{i\phi}v_*dz$. The derivative bound instead uses the three intervals
from zero to $7/8$, from $7/8$ to $9/8$, and from $9/8$ onwards.
It yields
\[
 \frac{|R'|}{v_0}\le\frac p{2\pi}
 \left(\frac{2}{pA(1-(7/8)^{p-1})}+\frac14+
       \frac{2}{pA((9/8)^{p-1}-1)}\right)<0.09p .
\]
Thus $R>p(1/2-0.8/(2\pi))>0$ and $R'/R<v_0/4$.

To bound $I'/I$, use $t=Tz^p$ on the positive real ray, for which
$xt^a=pTz$. Let $u=pTz$ and $\eta=2\pi a$. Apart from the common
factor $p e^{-Tz^p}/(2\pi)$ and the integration in $z$, the integrands
of $I$ and $I'/v_0$ are
\[
 \frac{e^u-e^{u\cos\eta}\cos(u\sin\eta)}z,\qquad
 e^u-e^{u\cos\eta}\cos(u\sin\eta+\eta).
\]
Both are nonnegative. If $z\ge3/8$ then
$e^{-u(1-\cos\eta)}\le e^{-12}$ as $\cos\eta<1/3$.
The second integrand is at least
$(3/8)(1-e^{-12})/(1+e^{-12})$ times the first.
The portion of $I$ below $3/8$ is less than $1/100$ of its total.
Indeed, omitting the common constant $p/(2\pi)$, that portion
is at most $(3/4)pT e^{(3/8)pT}$ since
$|e^u-e^{u e^{i\eta}}|\le2u e^u$.
The portion on $[9/10,1]$ is at least
$0.09e^{(0.9p-1)T}$. Their ratio is at most
$(0.75p/0.09)T e^{-(0.525p-1)T}<0.01$ for $T\ge12$ and
$4\le p\le5$. Therefore $I'/I>0.3v_0>R'/R$. This proves strict
phase increase on the unbounded range.

\paragraph{A compact uniform certificate.}
For $0\le x\le37$ and $1/5\le a\le1/4$, the rotated contour at
$S=128$ gives
\[
 |R_{xx}|\le8,\quad |R_{xxx}|\le16,\quad
 |R_{aa}|\le2\cdot10^6,\quad |R_{xaa}|\le6\cdot10^6 .
\]
Here are details for these deliberately loose constants.
Set $H=e^{2\pi ia}t^a$, $L=2\pi i+\log t$.
In the unsubtracted integral the four derivative multipliers are
$H^2,H^3,(xH+x^2H^2)L^2$, and
$(H+3xH^2+x^2H^3)L^2$.
On the vertical segment $|L|^2=(\log u)^2+25\pi^2/4$.
Thus the integral of $|H|^m|L|^2$ with measure
$|e^{-t+xH}dt/t|/(2\pi)$ is at most
\[
 \frac1{2\pi}\left[\frac{2}{(m/5)^3}
   +\frac{62}{m/5}+\frac{89\,128^{m/4}}{m/5}\right]
\]
on this segment, by splitting at $u=1$. Without $|L|^2$ its bound is
$128^{m/4}/(2\pi m/5)$. On the horizontal part $|t|^{ma-1}\le1$
for $1\le m\le3$, $|L|\le13+u/128$ and the damping is $e^{-u/2}$,
so the further contributions are at most $60$ with $|L|^2$ and
$1/\pi$ without it. These bounds imply the four constants.

The grid has step $1/8$ in $x$ and $1/2000$ in $a$. At its vertices
evaluate $R,R'$ through $n=383$ by 176-bit Arb arithmetic.
The error of bilinear interpolation of either function is at most
$(M_{xx}/8^2+M_{aa}/2000^2)/8$, where $M_{xx}$ and $M_{aa}$ bound
its respective pure second derivatives. This follows by applying the
one-dimensional linear interpolation error successively in the two
variables. On each order interval enclose the positive coefficients
$i_n=[x^n]I$ above and below. On a grid cell a lower bound for
$xI'/I$ is given by evaluating $\sum n i_n x^{n-1}$ from below at
the left endpoint and $\sum i_n x^{n-1}$ from above at the right.
This expression also has a positive limiting value at $x=0$.

A tail allowance $10^{-50}$ suffices for each truncated series involved.
In fact the four terms $n=4l,\ldots,4l+3$ with $l\ge96$ contribute
at most $4(4l+4)(l+1)!37^{4l+3}/(4l)!$, including a factor $n$
when necessary. The ratio of successive majorants is below $1/2$,
using $(l+2)/(4l+1)^4$ decreasing there. The first majorant is below
$5\cdot10^{-51}$. The verifier \texttt{kanter\_verify.py} checks
$R>0$ and $(x I'/I)R>xR'$ for all 29,600 closed grid cells by these
bounds. For $x>0$ this implies $(I/R)'>0$, completing the proof.

\paragraph{Consequences for stable laws.}
In the notation of the source, the one-sided free stable law
$X_\beta$ is a classical GGC for every $0<\beta\le4/5$. Indeed,
(3) expresses it as the independent product of $K_\beta$ and
$U^{1-1/\beta}$ with $U$ uniform on $(0,1)$. The GGC implication
for that product is precisely the reduction in the source's
Section~2.3.1. Its Section~2.3.2 then gives the extended Thorin
property for the two-sided $X_{\beta,\rho}$, $0\le\rho\le1$, throughout
the same order range. These are direct consequences of
Theorem~\ref{th:kanter}. For the ordinary positive stable law,
equation (12) and Remark~4(a) of the source also identify $[1/5,1)$
as the exact range on which the survival function of
$Z_a^{-a/(1-a)}$ is hyperbolically completely monotone.

\section{A shifted free stable law that is not a gamma mixture}
\label{sec:shifted-stable}
For $0<\beta<1$ let $X_\beta$ be the positive free stable variable
in the normalization of Hasebe, Simon and Wang, and let
$b_\beta=\inf\operatorname{supp}X_\beta$.
In Remarks~2(b) and~10(b) of \emph{Some properties of the free stable
distributions} the authors predict that $Y_\beta=X_\beta-b_\beta$
is a scale mixture of gamma distributions of shape $3/2$ for all
$0<\beta<1$. Both remarks concern the same conjecture.
\begin{theorem}\label{th:shiftedstable}
$Y_\beta$ is a gamma scale mixture with shape $3/2$ for
$0<\beta\le4/5$ and is not such a mixture for $23/25\le\beta<1$.
In particular the proposed all-order-parameter assertion is false.
\end{theorem}
The intermediate interval $(4/5,23/25)$ is not classified here.
For the positive range, the phase limit in Proposition~5 of the
source is $1/2$. With Theorem~\ref{th:kanter}, the argument of its
Remark~10(b) applies: the Thorin mass of $K_\beta$ is $1/2$, so
$K_\beta-b_\beta$ is a gamma scale mixture of shape $1/2$.
The same remark proves the implication from this last property to the
required shape-$3/2$ mixture of $X_\beta-b_\beta$.

\begin{proof}[Proof of failure near order one]
If the gamma mixture representation held, then $g_\beta(x)=
x^{-1/2} f_{Y_\beta}(x)$ would be completely monotone. In fact mixing
the gamma rates averages positive multiples of $e^{-tx}$ after division
by $x^{1/2}$. Proposition~6 of the source gives an all-orders expansion
$g_\beta(x)\sim\sum_{n\ge0}a_n(\beta)x^n$ at $0+$. Its formula implies
that $(-1)^n a_n(\beta)$ has the sign of
\[
 A_n(\beta)=[v^{2n}](1-v)^{-2}
   {}_2F_1(1+\beta,1;3;v)^{-n-1/2}.
\]
If $g_\beta$ is CM then $(-1)^n a_n(\beta)\ge0$ for all $n$.
Indeed $g_\beta$ has a finite limit at zero. The $n$th forward
difference at zero of each $e^{-t x}$, multiplied by $(-1)^n$, is
nonnegative. By the positive Laplace representation the same is true
for $g_\beta$. Dividing by the $n$th power of the mesh and taking
the mesh to zero gives $(-1)^n n!a_n(\beta)\ge0$ directly from
the expansion. Termwise differentiation of an asymptotic series is
unnecessary.

For the eight closed intervals $[i/100,(i+1)/100]$ with $i=92$ to
$99$, respectively, the following orders certify $A_n(\beta)<0$:
\[
  (71,55,45,35,30,29,24,23).
\]
Here is the finite rational calculation. Set $\beta=(i+u)/100$ and
$N=2n$, $d=-n-1/2$. The hypergeometric coefficients are
$b_0=1$ and $b_k=b_{k-1}(\beta+k)/(k+2)$. For the coefficients
of its $d$th power put $h_0=1$ and use logarithmic differentiation:
\[
 h_k=\frac1k\sum_{l=1}^k ((d+1)l-k)b_l h_{k-l},\qquad
 A_n=\sum_{k=0}^N(N+1-k)h_k.
\]
Write $A_n((i+u)/100)=\sum_{l=0}^N d_l u^l$. Its degree-$N$
Bernstein coefficients on $[0,1]$ are
$\sum_{l=0}^j d_l\binom jl/\binom Nl$, $0\le j\le N$.
The verifier \texttt{mixture\_rational.py} checks that every one
is strictly negative in each row. The Bernstein basis is nonnegative
and sums to one. This proves strict negativity on the entire union of
intervals and contradicts the necessary sign for the original mixture.
\end{proof}
The mixture obstruction tests the density at the left edge. The GGC
criterion in Section~\ref{sec:kanter} instead uses the phase of a
Laplace transform of a different variable. The two conditions lead
to separate conclusions in the source's terminology.

\section{Shifted gamma saddles and the stable-edge coefficient identity}
\label{sec:stirling-edge}
In Remark~12 of \emph{Some properties of the free stable distributions}
(Hasebe, Simon and Wang, \href{https://arxiv.org/abs/1805.01133}{arXiv:1805.01133}),
the authors conjecture an all-orders coefficient formula for the density
of their endpoint variable $Y_1$. Proposition~7 identifies its
coefficient $a_n(1)$ with
\[
 \frac{(-1)^n2^{2n+1/2}n!}{\pi(2n+1)!}(c_0+\cdots+c_n),
\]
where $c_n$ are the Stirling coefficients of the gamma function.
Their conjecture is that the following $b_n$ equals this partial sum:
\[
 b_n=\frac{1}{2^{2n+1/2}n!}\left.\frac{d^{2n}}{dv^{2n}}
 \left\{\frac1{(1-v)^2}
 \left(\frac{v^2}{-v-\log(1-v)}\right)^{n+1/2}\right\}\right|_{v=0}.
\]
All analytic fractional powers here have their positive value at zero.
We derive the identity from the shifted gamma integral, keeping a
real shift $d$ as a parameter.

\begin{proposition}
For $n\ge0$ put
\[
 B_n(d)=(2n-1)!![v^{2n}](1-v)^{-d-1}
 \left(\frac{v^2}{2(-v-\log(1-v))}\right)^{n+1/2},
 \qquad (-1)!!=1 .
\]
These are polynomials of degree exactly $2n$ in $d$, with $B_n(0)=c_n$.
They satisfy $B_0(d)=1$ and
$B_n(d+1)-B_n(d)=(d+1)B_{n-1}(d+1)$ for $n\ge1$.
In particular $b_n=B_n(1)=\sum_{j=0}^n c_j$ and hence the proposed
stable-edge formula holds in every order.
\end{proposition}
\begin{proof}
For $x>\max(d,0)$ apply the substitution $t=x(1-v)$
in the integral for $\Gamma(x-d)$. The normalized integral becomes
\[
 G_d(x)=\sqrt{\frac{x}{2\pi}}\,\frac{e^x\Gamma(x-d)}{x^{x-d}}
 =\sqrt{\frac{x}{2\pi}}\int_{-\infty}^{1}
 (1-v)^{-d-1} e^{-x\varphi(v)}\,dv,\qquad
 \varphi(v)=-v-\log(1-v).
\]
The phase is strictly positive away from its unique minimum at zero,
where $\varphi(v)=v^2/2+O(v^3)$. Contributions outside a fixed
neighborhood of zero are exponentially small: on that complement
bound $e^{-x\varphi}$ by $e^{-(x-x_0)\inf\varphi}e^{-x_0\varphi}$
for a fixed $x_0>\max(d,0)$. In that neighborhood use the analytic
coordinate $w=v\sqrt{2\varphi(v)/v^2}$, with inverse $v(w)$.
Expand $A(w)=(1-v(w))^{-d-1}v'(w)$ and integrate symmetrically.
The Gaussian moments give $G_d(x)\sim\sum_{n\ge0}x^{-n}
(2n-1)!![w^{2n}]A(w)$; the Taylor remainder through odd order
$2N+1$ is $O(x^{-N-1})$. By the change of variable in a
coefficient residue,
\[
 [w^{2n}]A(w)=[v^{2n}](1-v)^{-d-1}
 \bigl(v/w(v)\bigr)^{2n+1}.
\]
This yields the displayed coefficient, whose term of degree $2n$ in
$d$ comes only from $[v^{2n}](1-v)^{-d-1}$ and has nonzero multiplier.
At $d=0$ the normalized integral defines the ordinary Stirling coefficients.
Finally $G_{d+1}(x)=G_d(x)/(1-(d+1)/x)$ by the gamma recurrence.
Comparison of coefficients proves the finite-difference relation. At
$d=1$ the factorial conversions in the proposition give $B_n(1)=b_n$.
Putting $d=0$ in the recurrence gives the claimed partial sums.
\end{proof}
This derivation uses classical saddle-point analysis and coefficient
inversion. It avoids interchanging the limiting stable law with its
asymptotic expansion. Generalized Stirling coefficients with a shift
have long been studied, which limits the novelty attributable to the
adjacent coefficient identity.

\section{A sharp parameter half-plane for Jacobi endpoint comparison}
\label{sec:jacobi}
Booker, Str\"ombergsson and Then, \emph{Bounds and algorithms for the
$K$-Bessel function of imaginary order}, LMS J. Comput. Math. 16 (2013),
78--108, \href{https://doi.org/10.1112/S1461157013000028}{Conjecture~1,
equation (34)}, consider $f_n(x,s)=n!P_n^{(-s,s)}(x)$.
They establish the endpoint bound for $j=0$, for $r=0$, and for
$n\le j+100$. The following classification in asymmetric weights
includes all the derivative orders in their question.

\begin{theorem}\label{th:jacobi}
Fix $a,b>-1$. For nonnegative integers $m,j,k$ put
$Q(x,s)=P_m^{(a-s,b+s)}(x)$. The inequalities
\[
 |\partial_s^j\partial_x^kQ(x,s)|
       \le |\partial_s^j\partial_x^kQ(1,s)|
 \qquad(-1\le x\le1,\quad \Re s\le0)
\]
hold simultaneously for all $m,j,k$ if and only if
$a\ge b$ and $a\ge-1/2$. In particular the 2013 conjecture holds
at every $s=ir$, with no restriction on $n-j$.
\end{theorem}
\begin{proof}
First take $s=ir$, $j=k=0$, and $m\ge1$.
Write $y(x)=Q(x,ir)$ and $L=m(m+a+b+1)>0$.
The Jacobi equation yields the exact energy identity
\[
 \frac{d}{dx}\{L|y|^2+(1-x^2)|y'|^2\}
       =2[(a+b+1)x+a-b]|y'|^2.
\]
Suppose $a\ge b$ and $a\ge-1/2$. When $a+b+1\ge0$ the bracket
is nondecreasing, so the maximum of the energy is at an endpoint.
When $a+b+1<0$, its minimum is at $x=1$, with value $2a+1\ge0$,
and the same conclusion holds. The modulus of the endpoint at $+1$
is at least that at $-1$, since their values up to sign are
\[
 \frac1{m!}\prod_{h=1}^m(h+a-ir),\qquad
 \frac1{m!}\prod_{h=1}^m(h+b+ir)
\]
and $h+a\ge h+b>0$.

For a fixed $x$ now view the endpoint and interior polynomials as
functions of $s$, and write $P(s)=Q(1,s)$ and $H(s)=Q(x,s)$.
The zeros of $P$ lie strictly in the right half-plane.
The Jacobi binomial formula gives the same degree-$m$ leading
coefficient $(-1)^m/m!$ for $P$ and $H$ (the two binomial factors
contribute a binomial expansion of a constant).
The preceding boundary estimate and the maximum modulus principle
therefore give $|H/P|\le1$ throughout the closed left half-plane,
using $H/P\to1$ at infinity. For every $|u|<1$, the polynomial
$P+uH$ has degree $m$ and all zeros in the open right half-plane.
By Gauss--Lucas the same holds for each nonconstant derivative,
and its $m$th derivative is nonzero. For any $s$ in the closed left
half-plane, $|H^{(j)}(s)|>|P^{(j)}(s)|$ would contradict this on
taking $u=-P^{(j)}(s)/H^{(j)}(s)$. Higher derivatives vanish.
This elementary use of polynomial comparison explains the transfer
from the known zeroth derivative to all orders. For the $x$-derivatives
use the Jacobi differentiation formula
\[
 \partial_x^kQ(x,s)=2^{-k}(m+a+b+1)_k
       P_{m-k}^{(a+k-s,b+k+s)}(x),\qquad 0\le k\le m.
\]
Its prefactor is independent of $s$ and the shifted parameters
satisfy the same sufficient conditions. The constant case is immediate.

For necessity, $a<b$ fails already for $m=1$, $x=-1$, $s=0$.
If $b\le a<-1/2$, the bracket in the energy derivative is strictly
negative on some interval $(c,1]$. At $s=0$, for sufficiently large
$m$ the derivative of $P_m^{(a,b)}$ has a zero $x_m$ in $(c,1)$.
Indeed, it is a nonzero multiple of the degree-$m-1$ orthogonal
Jacobi polynomial for weight $(1-x)^{a+1}(1+x)^{b+1}$. Its largest
zero tends to $1$: the largest eigenvalue of the compression of
multiplication by $x$ to polynomials of degree at most $m-2$ is
that zero, by the three-term recurrence. Its Rayleigh quotient on
$(1+x)^{m-2}$ tends to $1$ by concentration of the weighted integrals
at the right endpoint. Integrating the strict energy decrease from
$x_m$ to $1$ gives $|Q(x_m,0)|>|Q(1,0)|$, since the derivative
term vanishes at both ends of this integration. This proves necessity.
\end{proof}
For $a=b=0$, the hypothesis on Conjecture~1 in the source's
Proposition~4 is consequently available for every positive integer
$n$. The energy method and Gauss--Lucas comparison are classical.
The full classification gives the precise obstruction to the same
endpoint choice for complex-shift Jacobi weights.

\section{The three global Bessel bounds in a 2009 conjecture}
\label{sec:red}
Baricz and Wu, \emph{Sharp exponential Redheffer-type inequalities for
Bessel functions}, Publ. Math. Debrecen 74 (2009), 257--278,
\href{https://doi.org/10.5486/PMD.2009.4291}{concluding remark 6 of
Section~1}, conjectured full ranges for three inequalities. We retain
their common source count. Let $j=j_{\nu,1}$ be the first positive
zero of $J_\nu$, let $y=x^2/j^2$, $R=(1-y)/(1+y)$,
$F=\mathcal J_\nu(x)=2^\nu\Gamma(\nu+1)x^{-\nu}J_\nu(x)$ (by its
even series for real $x$), and write $F_+=\mathcal J_{\nu+1}(x)$.
Their proposals are $F\ge R$ for $-1<\nu<\nu_0$, and $F\le R^\beta$,
$F_+/F\ge R^{-\gamma}$ for $\nu>-1$, throughout $0\le y<1$.
Here $\beta=j^2/[8(\nu+1)]$, $\gamma=\beta/(\nu+2)$, and
$\nu_0\in(1,2)$ is the unique order with $\beta=1$.
\begin{theorem}\label{th:red}
The first inequality holds throughout the interval exactly when
$-1<\nu\le1/2$. The second holds throughout exactly when $\beta\le1$.
The third holds throughout for $-1<\nu\le1/2$, but fails for every
sufficiently large $\nu$, even at the fixed relative position
$y=255/256$. Thus all three proposed full ranges are false. The
large-order failures of (1.2) and (1.3) also concern the ranges
printed in Theorem~1 of the original article.
\end{theorem}
\begin{proof}
Set $r_k=j^2/j_{\nu,k}^2$, so that $r_1=1$, and
$S=\sum r_k=j^2/[4(\nu+1)]=2\beta$ by the first Rayleigh sum.
The product formula gives $F=\prod_{k\ge1}(1-r_k y)$.
For the first inequality the log difference after removing $1-y$ is
\[
 D(y)=\log(1+y)+\sum_{k\ge2}\log(1-r_k y).
\]
It is strictly concave through $[0,1]$, and $D(0)=0$. Consequently
nonnegativity throughout is equivalent to
$2\prod_{k\ge2}(1-r_k)\ge1$. Each $r_k$, $k>1$, increases strictly
with $\nu$. In fact the Watson order derivative at any positive zero
$c=j_{\nu,k}$ is
\[
 \frac{d}{d\nu}\log c
 =2\int_0^\infty K_0(2c\sinh v)e^{-2\nu v}\,dv
\]
(\href{https://dlmf.nist.gov/10.21.E17}{DLMF 10.21.17}), a strictly
decreasing function of $c$ for fixed $\nu$. At $\nu=1/2$ one has
$r_k=1/k^2$ and $\prod_{k\ge2}(1-1/k^2)=1/2$ by telescoping.
The convergent positive product is strictly decreasing in $\nu$, by
factorwise comparison, with a strict inequality for $k=2$. This gives
the first range.

For the second, $F/(1-y)$ tends to a finite positive number as
$y\uparrow1$, ruling out $\beta>1$. If $S\le2$, then
\[
 -\log F-S\operatorname{artanh}y
 \ge-\log(1-y)+(S-1)y-S\operatorname{artanh}y
 \ge y-\log(1+y)\ge0 .
\]
The middle expression decreases in $S$ and was evaluated at $S=2$.

For the third, logarithmic differentiation of the product in $x$
and the Bessel recurrence give
\[
 F_+/F=\sum_{k\ge1}p_k(1-r_k y)^{-1},\qquad p_k=r_k/S.
\]
Here $M=\sum p_k r_k=j^2/[4(\nu+1)(\nu+2)]=2\gamma$, by the second
Rayleigh sum or the linear coefficient of the quotient series in $y$.
Whenever $\sum r_k^2\le2$ we have $M\le2p_1$. Jensen's inequality
and the same scalar estimate as above give
\[
 \log(F_+/F)-M\operatorname{artanh}y
 \ge-p_1\log(1-y)+(M-p_1)y-M\operatorname{artanh}y\ge0.
\]
For $\nu\le1/2$ one has $r_k\le1/k^2$ and $\sum r_k^2<2$.

To see the opposite behavior, take $\nu$ large, $h=\nu+1/2$.
On $0<x<h$ before the first zero let $u=J_{\nu+1}(x)/J_\nu(x)$.
It satisfies $u'=1-2hu/x+u^2$. The smaller root of this right-hand
side is $b(x)=x/(h+\sqrt{h^2-x^2})$, a strictly increasing function.
Near zero, $u\sim x/(2\nu+2)< b$. A first contact from below would
give $u'=0<b'$, an impossibility. Thus $F_+/F\le
2(\nu+1)/(h+\sqrt{h^2-x^2})$. The standard first-zero asymptotic
$j_{\nu,1}/\nu\to1$ (\href{https://dlmf.nist.gov/10.21}{DLMF,
Section~10.21(vii)}) ensures $x=j\sqrt{255/256}<h$ eventually.
The upper limit of $F_+/F$ is at most $2/(1+1/16)<2$ there,
whereas $R^{-\gamma}\to511^{1/8}>2$. This proves the refutation.
\end{proof}

\section{Fractional degrees of two trigamma remainders}\label{sec:degrees}
F. Qi, \emph{Completely monotonic degree of a function involving trigamma
and tetragamma functions}, AIMS Mathematics 5 (2020), 3391--3407,
\href{https://doi.org/10.3934/math.2020219}{Remark~7.4},
conjectures degree four for each of two signed remainders. The same
assertion appeared in the earlier \href{https://arxiv.org/abs/1301.0154}{preprint},
equation (4.8) in version 3. Put
\[
 \Psi(x)=\psi'(x)^2+\psi''(x),\qquad
 h_\lambda(x)=\Psi(x)-\frac{x^2+\lambda x+12}{12x^4(x+1)^2},\quad x>0 .
\]
For the positive functions vanishing at infinity in this section, the
degree is the supremum of the real $r$ for which $x^rf(x)$ is completely
monotone (CM), meaning $(-1)^j(x^rf(x))^{(j)}\ge0$ for every $j\ge0$.
The source establishes CM of $h_\lambda$ for $\lambda\le0$ and of
$-h_\mu$ for $\mu\ge4$.

\begin{theorem}\label{th:trigamma-degrees}
For every $\mu\ge4$ the degree of $-h_\mu$ equals $3$.
The degree of $h_0$ equals $2$. There is a continuous strictly increasing
bijection $C:(0,1)\to(0,\infty)$, given below by an integral maximum,
such that for $\lambda<0$ the degree of $h_\lambda$ equals
$2+C^{-1}(-\lambda/12)$. All these degrees are attained.
\end{theorem}

Already the recurrence for the polygammas at zero gives
\[
 h_\lambda(x)=-\frac{\lambda}{12x^3}
 +\left(2\zeta(2)-\frac{37}{12}+\frac{\lambda}{6}\right)x^{-2}
 +O(x^{-1}).
\]
A nonzero nonnegative decreasing function cannot tend to zero at the
left endpoint. These expansions alone contradict the two proposed
degrees of four. The exact replacement requires controlling the signs
of fractional derivatives over a whole half-line.

Let $b(t)=t/(1-e^{-t})$, $K=b*b-tb$ with additive convolution, and define
\begin{align*}
 V(t)&=K''(t)-(t-2)-(2+13t/12)e^{-t},\\
 W(t)&=1-(11+9t)e^{-t}/12-K'''(t),\qquad
 j(t)=1-(1+t)e^{-t}.
\end{align*}
The trigamma transform $\psi'=\mathcal L b$ gives $\Psi=\mathcal L K$,
and $K=O(t^3)$ at zero. Partial fractions and Laplace differentiation yield
\[
 x^2h_\lambda=\mathcal L(V-\lambda j/12),\qquad
 -x^3h_\mu=\mathcal L\bigl(W+(\mu-4)te^{-t}/12\bigr).
\]

\begin{lemma}\label{lem:degree-kernels}
$V(t),W(t)>0$ for every $t>0$. At zero $V(t)=t^2/6+O(t^3)$.
At infinity $V(t)\sim t^3e^{-t}/6$ and $V'(t)\sim-t^3e^{-t}/6$.
\end{lemma}
\begin{proof}
The convolution identity
\[
 (1-e^{-s})^{-1}(1-e^{-(t-s)})^{-1}
 =\frac{(1-e^{-s})^{-1}+(1-e^{-(t-s)})^{-1}-1}{1-e^{-t}}
\]
gives the convergent expansion $K(t)=\sum_{k\ge0}P_k(t)e^{-kt}$, where
\[
 P_k(t)=t^3/6-t^2+2(\zeta(2)+H_k^{(2)})t
                -4(\zeta(3)-H_k^{(3)}),\qquad
 H_k^{(r)}=\sum_{l=1}^k l^{-r}.
\]
Indeed, in the convolution the remaining integral in the numerator is
$2\int_0^t s(t-s)/(e^s-1)\,ds$. Expand the denominator into exponentials
and integrate to obtain the formula. The series and every fixed
derivative converge absolutely locally on $t>0$. Its first exponential
implies the stated large-$t$ asymptotics.

Here are reproducible sign bounds. For $t\ge32$ one has
$P_k\ge t^3/8$, $0<P'_k\le t^2/2+8$, $P''_k=t-2$ and $P'''_k=1$.
After the $k=0$ cancellations, the $k$th polynomial contributing to
$V$ is $P''_k-2kP'_k+k^2P_k$. This is at least
$k^2(t^3/8-t^2-16)$ for $k\ge1$. For $W$ the polynomial is
$k^3P_k-3k^2P'_k+3kP''_k-1$, which is at least
$k^3(t^3/8-3t^2/2-25)$. The $k=1$ terms alone exceed the respective
additional subtractions $(2+13t/12)e^{-t}$ and $(11+9t)e^{-t}/12$.

On $2\le t\le32$ apply the same expansion to $e^t V$ and $e^t W$,
retaining $k\le32$. The absolute remainder of either is bounded by
$2\cdot10^5\,33^3/7^{32}$. Indeed each polynomial in the preceding
display is bounded absolutely by $10^5 k^3$ on this interval, and
$e^{-t}<1/7$. On $0\le t\le2$ compute the Taylor polynomials of
$V,W$ through degree $80$ from
\[
 [t^{m+1}]K=\sum_{l=0}^m b_l b_{m-l}\frac{l!(m-l)!}{(m+1)!}-b_m,
 \qquad b_l=[t^l]b.
\]
They have zero constant and linear terms, and $[t^2]V=1/6$. On
$|z|\le4$ one has $|b(z)|<7$ by the even Bernoulli coefficient formula
$|b_{2l}|=2\zeta(2l)/(2\pi)^{2l}$ for $l\ge1$. Integration along the
straight segment gives $|K(z)|\le224$ there. On $|z|\le3$ the Cauchy
derivative bounds for $K''$ and $K'''$ and the explicit exponential
terms bound $|V|,|W|$ by $2000$. Hence the discarded tails after
division by $t^2$ are at most $1500(2/3)^{81}$ on $[0,2]$.
The verifier \texttt{trigamma\_degrees.py} evaluates the two polynomial
bounds and the finite exponential sums by outward-rounded interval
arithmetic (100 bits). Bisection accepts respectively 11 and 166 boxes.
\end{proof}

\paragraph{An auxiliary transfer for fractional orders.}
For $0\le s<1$ define
\[
 (L_s u)(t)=\frac1{\Gamma(1-s)}\int_0^t(t-v)^{-s}u(v)\,dv.
\]
For $0<s<1$ set
\[
 C(s)=\max_{t>0}\frac{-L_s V'(t)}{L_s j'(t)} .
\]
The denominator is positive since $j'(t)=te^{-t}$. Notice that this is an
explicit scalar optimization with a strictly invertible dependence on the
degree. It involves no searches over derivative orders. Indeed, since
$V(0)=j(0)=0$, the Laplace transform of $L_s(V'+c j')$ is exactly
$x^{2+s}h_{-12c}(x)$. By Bernstein's theorem and uniqueness of Laplace
transforms this function is CM if and only if $c\ge C(s)$.

We verify all the asserted properties of the maximum. Write $R_s(t)$
for the quotient in its definition. As $t\downarrow0$, the expansions
$V'(t)=t/3+O(t^2)$ and $j'(t)=t+O(t^2)$ show $R_s(t)\to-1/3$.
As $t\to\infty$, uniformly for $0\le s\le s_0<1$,
\[
 R_s(t)=\frac{s\int_0^\infty V(v)\,dv}{t}+O_{s_0}(t^{-2}).
\]
To see this, split each integral at $t/2$. On the first part expand
$(1-v/t)^{-s}=1+sv/t+O_{s_0}(v^2/t^2)$. On the second part use the
exponential decay of $V',j'$ and the integrability of $(t-v)^{-s}$,
uniform for $s\le s_0$. Now $\int V'=0$, $\int vV'=-\int V$, and
$\int j'=1$. This proves the expansion. Thus for $s>0$ the maximum
is positive, finite, and attained at some interior $t$. On compact
$s$-subintervals of $(0,1)$ the two endpoint bounds are uniform and
the quotient is jointly continuous on compact interior rectangles.
Consequently $C$ is continuous there.

If $0<s_1<s_2<1$, the fractional integral semigroup identity makes
$R_{s_1}(t)$ a weighted average of $R_{s_2}(v)$, $0<v<t$, with strictly
positive weight $(t-v)^{s_2-s_1-1}L_{s_2}j'(v)$. The identity follows
by interchanging integrals and using the beta integral. Evaluate at
a maximizer for $s_1$. Since the averaged quotient is negative near
$v=0$ but $C(s_2)>0$, strict inequality $C(s_1)<C(s_2)$ follows.
When $s\downarrow0$, the compact-interior limit is $R_0=-V/j<0$,
and the uniform endpoint estimates give $C(s)\to0$. When $s\uparrow1$,
$L_su(t)\to u(t)$ at fixed $t>0$ for continuous $u$, by concentration
of the normalized power kernel. Therefore
$\liminf_{s\uparrow1} C(s)\ge-V'(t)/j'(t)$ for every $t$. The latter
quotient grows as $t^2/6$ by Lemma~\ref{lem:degree-kernels}. Hence
$C(s)\to\infty$.

Finally Lemma~\ref{lem:degree-kernels} proves the two integer degrees
are attainable, and the pole orders at zero give their upper bounds.
The strictly increasing inverse gives the exact attained degree when
$\lambda<0$, since CM at one power implies CM at every smaller power
by multiplication by a CM power. This proves
Theorem~\ref{th:trigamma-degrees}.

\section{Reproducing the finite checks}
Everything needed to typeset this article is in this file. Compile with
\texttt{pdflatex} twice, or until all cross-references stabilize.
After \verb|\end{document}| the source contains five files between
matching \verb|%% BEGIN| and \verb|%% END| lines. Copy only the lines
between a matching pair of delimiters.
\begin{itemize}
\item Run \texttt{roy\_square\_verify.py} with Python~3. It uses exact
integer arithmetic. Its finite differences certify the sign of infinitely
many coefficients on arithmetic progressions.
\item Build \texttt{zero\_forcing\_audit.cpp} using a C++20 compiler
and run the resulting binary.
It simulates the coloring rule on the 34-vertex example.
\item Run \texttt{mixture\_rational.py} with Python-FLINT. It computes
the Bernstein coefficients of eight polynomials over the rationals
and checks the sign of each coefficient.
\item Run \texttt{kanter\_verify.py} and \texttt{trigamma\_degrees.py}
with Python~3 and \texttt{python-flint} 0.9. The displayed final
decimal margins are informational. Each required inequality is tested
using Arb balls or exact rationals before a success message is emitted.
\end{itemize}
The signs on unbounded domains require the analytic tail estimates
given in the proofs. The programs for compact sets do not replace
those estimates.
\appendix

\section{Zero-forcing polynomials of joins of clique unions}
\label{sec:zfjoin}
Boyer, Brimkov, English, Ferrero, Keller, Kirsch, Phillips and Reinhart,
\emph{The zero forcing polynomial of a graph}, Discrete Applied Mathematics
258 (2019), 35--48, \href{https://doi.org/10.1016/j.dam.2018.11.033}{doi:10.1016/j.dam.2018.11.033},
posed Conjecture~1 in Section~6 of the
\href{https://arxiv.org/abs/1801.08910}{2018 preprint}. It asserts
unimodality of the coefficients of $\mathcal Z(G;x)$ for every graph.
P. Graham, \emph{A graph whose zero-forcing polynomial is not unimodal},
1 November 2026, \href{https://doi.org/10.5281/zenodo.23875153}{Zenodo
23875153}, had already refuted this conjecture with a $73$-vertex graph.
That construction also uses a join of unions of cliques and the fact
that twin vertices cannot both stay white. We give a smaller witness,
a formula for uniform clique sizes on each side, and a construction
with arbitrarily many deep valleys. The earlier disproof is excluded
from the count of new source-conjecture resolutions.

In this polynomial $z(G;k)=[x^k]\mathcal Z(G;x)$ counts the initially
blue $k$-element sets for which repeated application of the following
rule colors every vertex: a blue vertex with exactly one white
neighbor may color that neighbor blue.

\begin{theorem}
Let $a,b,r,s\ge2$. Form $G=(a K_r)\vee(b K_s)$ by taking the join
of two disjoint unions of cliques and write $N=ar+bs$. For the polynomial
$F(u)=\sum_{t\ge0}z(G;N-t)u^t$ one has
\begin{align*}
F(u)={}&(1+ru)^a+(1+su)^b-1\\
 &+bsu\big[(1+ru)^a-1-(ru)^a\big]\\
 &+aru\big[(1+su)^b-1-(su)^b\big]-abrsu^2 .
\end{align*}
The connected $34$-vertex graph $(7K_2)\vee(4K_5)$ gives
\[
 z(G;28)=13888\ >\ z(G;29)=11872\ <\ z(G;30)=13785 .
\]
In particular the conjecture is false even on cographs of diameter two.
\end{theorem}
\begin{proof}
The complement of the initial blue set will be called the white set.
Two white vertices in the same constituent clique cannot both eventually
be colored: before the first of them changes color every blue vertex
adjacent to either is adjacent to both. Thus a forcing set leaves at most
one white vertex in each constituent clique. Write $i,j$ for the numbers
of initially white vertices on the two sides. If $i,j\ge2$ no first
move is available since all vertices are adjacent to the entire opposite
side.

If $j=0$, each left clique containing a white vertex contains a blue
vertex with exactly that one white neighbor. All of the whites can be
colored. Interchanging the sides handles $i=0$. These cases give the
first line of the formula. Suppose now that $j=1$ and $i\ge1$. If $i<a$,
a blue vertex in an entirely blue left clique can color the only white
vertex on the right. The previous argument then applies. If $i=a$, every
left blue vertex sees two whites and every right blue vertex sees at least
$a\ge2$ whites, so the process cannot start. The contribution for $j=1$
and $1\le i<a$ is the second line. The last line adds the analogous
case $i=1$ and $1\le j<b$, removing the duplicated $i=j=1$ contribution.
This exhausts all possible initial colorings. Substitution yields the
displayed three exact integers. A join of disjoint unions of cliques is
a cograph by the disjoint-union and join characterization. Since both
sides have at least two cliques, its diameter is two.
\end{proof}

\paragraph{Two concentrations.}
The same family can give a valley of arbitrarily large depth. Set
$r=2$, $s=8$, $a=2b$ and let $b\to\infty$. The two polynomial terms
$P=(1+2u)^{2b}$ and $Q=(1+8u)^b$ have the same total mass $9^b$.
Their normalized coefficients are binomial laws with respective means
$4b/3$ and $8b/9$, separated by $4b/9$. In the coefficientwise order
for nonnegative polynomials, the proof gives
\[
 P\preceq F,\qquad Q\preceq F,\qquad
 F\preceq (1+8bu)P+(1+4bu)Q .
\]
At the nearest integer to the midpoint $10b/9$, the coefficients of $P,Q$
and their unit shifts are all $9^b O(e^{-cb})$ for some absolute $c>0$
by Hoeffding's inequality. In each of the two windows of radius $b/9$
about the means, a coefficient of the corresponding pure polynomial is
at least $9^b/(5b)$ when $b$ is large, by the same concentration bound
and the number of integer indices in a window. Therefore the valley
coefficient divided by the smaller of these two peak coefficients is
$O(b^2e^{-cb})$. This argument concerns the original ordinary color-change
rule and the unweighted cardinality enumerator.

\paragraph{Any prescribed number of separated peaks.}
Fix $h\ge2$, and choose $M$ tending to infinity through multiples of
$\operatorname{lcm}(1,\ldots,h)$. Take the join of $h$ graphs
$G_j=a_j K_{r_j}$ with $a_j=M/j$ and $r_j=3^j-1$. The same first-move
obstruction shows that all white vertices lie in one side, or that they
lie in two sides with exactly one white vertex on one of these two sides.
In fact if three sides have white vertices, every blue vertex sees at
least two whites outside its side. Two sides with at least two whites
each give the same obstruction. There is still at most one white per
constituent clique. Conversely, the single-side case always forces
under this last restriction. For
$P_j(u)=(1+r_ju)^{a_j}$ and $N=\sum_j a_jr_j$, this gives
\[
 P_j\preceq F\quad(1\le j\le h),\qquad
 F\preceq (1+Nu)\sum_{j=1}^h P_j .
\]
All $P_j$ have mass $3^M$. Their means
$M(1-3^{-j})/j$ are pairwise distinct with minimum spacing
$\delta_h M$ for some $\delta_h>0$. Indeed
$(1-3^{-j})/j$ is $2/3$ times the mean of the strictly decreasing
list $1,1/3,\dots,3^{1-j}$. For each fixed $h$ these binomial
distributions put all but $O_h(e^{-c_hM})$ of their mass within
distance $\delta_h M/4$ of their respective means, by Hoeffding's
inequality. At each intervening midpoint all the coefficients
of all $P_j$ and their unit shifts are $O_h(3^M e^{-c_h'M})$.
Each mean window contains a coefficient at least
$3^M/(2(N+1))$. Since $N=O_h(M)$, the graph has $h$ separated
peaks whose intervening valleys are smaller by arbitrarily large
factors. The construction stays within connected cographs of diameter two.

\section{Finite derivative zero counts and infinite divisibility}
\label{sec:hsw-bell}
T. Hasebe, T. Simon and M. Wang, \emph{Some properties of the free
stable distributions}, Ann. Inst. H. Poincar\'e Probab. Statist. 56
(2020), 296--325, \href{https://arxiv.org/abs/1805.01133}{arXiv:1805.01133},
Remark~15(b), define $\mathrm{BS}_n$ by the requirement that the $k$th
derivative of a smooth density on $(0,\infty)$ have exactly
$\min(k,n)$ distinct zeros there for every $k\ge0$. They conjecture
that all such densities are classically infinitely divisible for every
$n\ge2$. The same remark identifies the class with the earlier
$\mathrm{WBS}_{n-1}$ of T. Simon, \emph{Positive stable densities and
the bell-shape}, \href{https://arxiv.org/abs/1302.0956}{arXiv:1302.0956}.
That definition specifies endpoint limits in addition to interior
zero counts. The two formulations have different answers.
The positive conclusion with those extra conditions was already
established by Theorem~1.13 and Corollary~1.4 of M.~Kwa\'snicki and
T.~Simon, \emph{Characterisation of the class of bell-shaped functions},
\href{https://arxiv.org/abs/1910.07752}{arXiv:1910.07752}. We give an
explicit counterexample to the separate zero-count definition and a
direct Gaussian-variation proof to pinpoint the role of the endpoint.

\begin{theorem}\label{th:hsw}
For every $0<a\le1/10$, the probability density proportional
to $e^{-x}x(x+a)^7$ on $x>0$ belongs to the printed class
$\mathrm{BS}_8$ and is not infinitely divisible.
Conversely, if $f\in\mathrm{BS}_n$ and
$f^{(j)}(0+)=0$ for $0\le j<n$, then its law is infinitely
divisible. The latter conclusion includes the earlier endpoint
convention in $\mathrm{WBS}_{n-1}$.
\end{theorem}
\begin{proof}
First consider a real monic polynomial $P$ of degree $n$ with all
roots in $(-\infty,0]$. Write $D$ for differentiation and $T=1-D$.
The $k$th derivative of $e^{-x}P(x)$ is $(-1)^ke^{-x}T^kP(x)$.
For any real-rooted monic polynomial $Q$ with distinct roots $r_i$
of multiplicities $m_i$, the roots of $TQ$ comprise $r_i$ with
multiplicities $m_i-1$ and the solutions of
$\sum_i m_i/(x-r_i)=1$. There is one simple solution in each gap
$(r_i,r_{i+1})$ and one to the right of the largest root.
Consequently each iteration increases the number of strictly positive
roots, counted with multiplicity, by at most one. Once all $n$ roots
are positive and simple, they stay positive and simple at every
subsequent iteration.

Apply this to $P_a=x(x+a)^7$ with $n=8$. At $a=0$, $T^kx^n$ has
a zero of multiplicity $n-k$ at zero for $k<n$ and exactly $k$
positive simple roots for $k\le n$, by the same gap argument
starting from $x^n$. By continuity of the coefficients, each
distinct positive simple root persists under sufficiently small
positive $a$. There are only finitely many orders $k\le8$ to check.
Together with the upper count and propagation for $k>8$, this
proves the required counts for \emph{all} derivatives. This count extends through $0<a\le1/10$.
Indeed, for $1\le k\le8$ the derivative at zero is the alternating
sum with leading positive term and unsigned terms, in descending $j$,
\[
 \frac{k!}{(k-j)!}\binom7{j-1}a^{8-j},\qquad k\ge j\ge1 .
\]
The ratio of each succeeding term to its predecessor is
$a(j-1)/[(9-j)(k-j+1)]\le7a<1$, so no such derivative has a root
at zero. In this parameter interval the positive root counts cannot
change: all roots remain real and bounded on compact parameter
intervals, and none can cross zero. Positive roots stay simple by
induction in $k$ and the gap rule. In particular the example vanishes
at the left endpoint although $f'(0+)>0$.

The unnormalized Laplace transform at $s$ is
$(1+s)^{-9}Q(a(1+s))$, by integrating the polynomial termwise, where
\[
 Q(y)=y^7+14y^6+126y^5+840y^4+4200y^3
                +15120y^2+35280y+40320 .
\]
This polynomial has a root $z$ with $\Re z>0$. Indeed its fifth
Hurwitz determinant, in the descending-coefficient convention, is
\[
 \det\begin{pmatrix}
 14&840&15120&40320&0\\
 1&126&4200&35280&0\\
 0&14&840&15120&40320\\
 0&1&126&4200&35280\\
 0&0&14&840&15120
 \end{pmatrix}=-423901900800 .
\]
Were all roots in the closed left half-plane, $Q(y+u)$ for every
$u>0$ would be strictly Hurwitz stable, so all its leading Hurwitz
determinants would be positive. Continuity as $u\downarrow0$
contradicts the display. In fact a root satisfies $\Re z>1/10$:
the same fifth minor for $Q(y+1/10)$ is
$-26971892946311754902293/61035156250<0$ by binomial expansion.
For $0<a<\Re z$, the Laplace transform
therefore has the zero $s=z/a-1$ in the open right half-plane.

For completeness no nonnegative infinitely divisible law can have such
a zero. In every $m$-fold equal convolution decomposition, the
summands are nonnegative almost surely, since the essential infimum
of the sum of independent identically distributed variables is $m$
times the common essential infimum. Their Laplace transforms are
holomorphic on the open right half-plane. The multiplicity of any
zero of the transform of the full law must thus be divisible by
every positive integer $m$, an impossibility. This proves the first
assertion.

\paragraph{Gaussian smoothing and endpoint derivatives.}
For the second assertion, extend the density by zero to $\mathbb R$.
Every derivative tends to zero at $+\infty$. In fact $f$ is eventually
monotone and integrable there. Inductively $f^{(j)}$ is eventually
monotone and any nonzero or infinite limit would contradict the
preceding derivative's vanishing limit. Because $f^{(j)}$ has constant
sign on finitely many intervals, the endpoint conditions imply
integrability of $f^{(j)}$ on the half-line for $1\le j\le n$.
This follows by integrating on each such interval and using the
endpoint limits of $f^{(j-1)}$. The distributional $j$th derivative
of the extension is the zero extension of $f^{(j)}$ for $j\le n$,
with no boundary atoms, by integration by parts.

Let $p$ be any centered Gaussian density of positive variance.
For $k\le n$ the Gaussian variation-diminishing property gives at
most $k$ sign changes of $(p*f)^{(k)}$, by applying it to $f^{(k)}$
extended by zero. For $k=n+d$, $d>0$, write
$h=f^{(n)}\mathbf1_{(0,\infty)}$ (for $n=0$ take $h=f$).
It is integrable, so $h_\epsilon=h\mathbf1_{[\epsilon,\infty)}$
converges in $L^1$ to $h$. Its distributional $d$th derivative
consists of $f^{(k)}$ on $(\epsilon,\infty)$ together with a linear
combination of point-mass derivatives at $\epsilon$ of orders at most
$d-1$. The regular part is an integrable signed measure there since
$f^{(k-1)}$ tends to zero and $f^{(k)}$ has finitely many zeros.
The point-supported part can be approximated by a signed measure on
the $d$ points $\epsilon-j\delta$, $1\le j\le d$, as
$\delta\downarrow0$. Indeed, match its action on polynomials through
degree $d-1$ by solving a Vandermonde system. The weights are
$O(\delta^{1-d})$ and the Taylor remainder then gives convergence
against smooth tests, including each translate of the Gaussian.

The approximating measures including the regular part have at most
$d+n=k$ sign changes. There are at most $d-1$ among the added atoms,
at most one at the join, and at most $n$ on the right.
Gaussian convolution cannot increase the number of sign changes.
This classical variation-diminishing theorem follows from the total
positivity of the Gaussian kernel; see I.~J. Schoenberg,
\emph{On P\'olya frequency functions. I}, J. Analyse Math. 1 (1951),
331--374. Passing first to the pointwise limit $\delta\downarrow0$
and then to $\epsilon\downarrow0$ shows that $p^{(d)}*h=
(p*f)^{(k)}$ also has at most $k$ sign changes. For the second limit
use boundedness of $p^{(d)}$ and the $L^1$ convergence. A strict
alternation at finitely many sample points persists under pointwise
convergence, which justifies both upper bounds at the limit.

All derivatives of $p*f$ tend to zero at both infinities by integrability
of $f$, and $p*f$ is positive. Successive mean-value arguments give at
least $k$ sign changes of its $k$th derivative: each sequence of
alternating nonzero values for the preceding derivative gives one more
such alternation for the derivative by taking differences between
successive values and using zero limits at both ends.
Thus all Gaussian smoothings are strictly bell-shaped in the sense of
M. Kwa\'snicki and T. Simon, \emph{Characterisation of the class of
bell-shaped functions}, \href{https://arxiv.org/abs/1910.07752}{arXiv:1910.07752}.
Their Corollary~1.4 implies infinite divisibility. Equivalently the
zero-extended original probability is weakly bell-shaped and that
corollary applies directly to it.
\end{proof}
The alternative proof of the known endpoint implication uses an
established characterisation after explicitly controlling the point-mass derivatives of a truncated half-line density.
The negative example addresses only the definition in Remark~15(b) of
the 2020 source. It is compatible with the authors' separate
gamma-mixture prediction, since the example divided by $x^8$ is
proportional to $e^{-x}(1+a/x)^7$, a completely monotone function.

\section{The gamma-mixture question for finite bell shape}
\label{sec:hsw-gamma}
Hasebe, Simon and Wang, Remark~15(b) cited in
Section~\ref{sec:hsw-bell}, also predict that every
$\mathrm{BS}_n$ density is a scale mixture of gamma densities of
shape $n+1$. This statement concerns a different property of the law.
The next theorem distinguishes the two notions of zero counts in this
question. To be a gamma mixture of this shape a probability density
$f$ must and need only have $x^{-n}f(x)$ completely monotone.
Indeed each positive-rate gamma density has the form
$v^{n+1}x^ne^{-vx}/n!$, and conversely normalization of a positive
Laplace representation gives a probability measure on those rates
by Tonelli's theorem. For the stronger endpoint convention,
Kwa\'snicki--Simon, Theorem~1.13, already give a convolution of $n$
exponentials and a completely monotone density. Conditional on the
mixing rate of the latter, its random variable is a sum of $n+1$
independent exponentials of possibly distinct rates. Writing unit
exponentials as their sum times their independent Dirichlet proportions
exhibits this as a gamma mixture of shape $n+1$. The result below
allows all lower endpoint jets to be nonzero.

\begin{theorem}\label{th:hsw-gamma}
If for every $j\ge0$ the derivative $f^{(j)}$ has exactly
$\min(j,n)$ distinct zeros in $(0,\infty)$ and changes sign at each
of those zeros, then $f$ is a gamma mixture of shape $n+1$.
Under the published zero-count definition without this sign-change
condition there is a counterexample already for $n=3$.
\end{theorem}
\begin{proof}
For the counterexample, take $f$ proportional to
$e^{-x}x((x-1)^2+6)$ on $(0,\infty)$. This is positive and has
$f(0+)=0$. Ignoring the fixed positive normalization, its first
two derivatives are
\[
 f'=-e^{-x}(x-1)((x-2)^2+3),\qquad
 f''=e^{-x}(x-2)(x-3)^2 .
\]
The gap argument for $1-D$ in Theorem~\ref{th:hsw} gives three
simple positive zeros for $f'''$ and every subsequent derivative.
Thus $f\in\mathrm{BS}_3$ in the printed sense. But its quotient by
$x^3$ is proportional to $e^{-x}(1-2/x+7/x^2)$. Its inverse Laplace
measure, up to that constant, is
\[
 \delta_1+\mathbf1_{v>1}\bigl[-2+7(v-1)\bigr]\,dv .
\]
This has negative mass on intervals immediately to the right of one.
Uniqueness of Laplace transforms rules out a positive representing
measure, so there is no gamma mixture.

For the positive assertion the tail sign of $f^{(j)}$ is $(-1)^j$.
Indeed all derivatives tend to zero at infinity as in the preceding
proof. Since $f>0$ and each derivative is eventually monotone and
nonzero, induction gives the alternating signs there. For $j\le n$
the $j$ sign changes force $f^{(j)}>0$ near zero. In particular
$f^{(r)}(0+)$ exists, is finite and is nonnegative for $0\le r<n$:
both this derivative and the following derivative are positive near
zero.

Fix a derivative order $k$ and set $m=n+k+2$.
Define the signed measure $\mu$ by pushing forward
$(-1)^m t^{m-1} f^{(m)}(t)\,dt/(m-1)!$ under $s=m/t$.
For $0\le r<n$ it is absolutely integrable against $s^r$ and
\[
 \int s^r\,d\mu(s)=(-1)^r
 \frac{m^r(m-1-r)!}{(m-1)!}\,f^{(r)}(0+).
\]
These formulas follow by successive integration by parts. Here are
endpoint details. If a derivative $f^{(r)}$ has finite limit at an
endpoint, then $t^j f^{(r+j)}(t)\to0$ there for every $j\ge1$.
For $j=1$ use eventual monotonicity of $f^{(r+1)}$ and secants on
$[t/2,t]$ and $[t,2t]$. Iterating these inequalities with the
previous weighted limits proves the assertion for every $j$.
These bounds hold at infinity and at zero for the $r$ in question.
Applying integration by parts first on compact intervals proves
the formula at the endpoints. Since $f^{(m)}$ has finitely many
zeros, convergence of these improper moments entails absolute
convergence. For $n=0$ the assertion follows immediately from
the alternating tail signs on the entire half-line; hence assume $n\ge1$.
Taylor's formula with integral remainder at a right endpoint tending
to infinity gives the exact identity
\[
 f(x)=\int (1-xs/m)_+^{m-1}\,d\mu(s).
\]
The Taylor polynomial at the right endpoint vanishes in that limit
by the same secant estimates.

For $x>0$ define
\[
 H(s)=(-\partial_x)^k\left[x^{-n}(1-xs/m)_+^{m-1}\right].
\]
For $0\le r\le n$ one has $(-1)^rH^{(r)}(s)\ge0$. In fact, applying
$(-\partial_s)^r$ first leaves a positive constant times
$x^{r-n}(1-xs/m)_+^{m-1-r}$. The first factor has alternating
derivatives of all orders, and the second has alternating derivatives
in $x$ through order $k$. Since $m=n+k+2$, these statements and the
differentiations are valid across the cutoff $xs=m$. Differentiating
under the measure integral is legitimate on a neighborhood of each
$x>0$: all nonzero terms have $s$ in a common compact interval.

The signs of $\mu$ as $s$ increases start positive and alternate at
$n$ points $0<b_1<\dots<b_n$. Interpolate $H$ at those points by
the polynomial $P$ of degree at most $n-1$. Its Newton divided
difference of order $j$ has sign $(-1)^j$ or is zero, by the derivative
bound. Each Newton basis product of $j$ factors $s-b_i$ has
coefficient of degree $r$ of sign $(-1)^{j-r}$. It follows that
$(-1)^r[s^r]P\ge0$ at every degree $r$. The moment formulas imply
$\int P\,d\mu\ge0$.
The interpolation remainder has sign
$(-1)^n\prod_i(s-b_i)$ or is zero, by the mean-value formula for
the divided difference of order $n$. This is precisely the sign
of $\mu$. Hence $\int(H-P)\,d\mu\ge0$ as well.
We have shown $(-\partial_x)^k(x^{-n}f(x))=\int H\,d\mu\ge0$,
which proves the gamma-mixture statement for every $k$.
\end{proof}
More generally the interpolation proof works whenever the first $n$
endpoint jets are finite and nonnegative and $f^{(m)}$ changes sign
at most $n$ times for each $m$; one uses the actual number
$\ell\le n$ of sign changes as the interpolation order, integrating
$H$ directly if $\ell=0$. The estimates at the endpoints in this
argument require the eventual monotonicity of each derivative,
which already follows from the assumed finite sign changes of the
following derivative. Thus no equality of the first $n$ endpoint
jets to zero is required. The source's shape-$n+1$ mixture assertion and its
infinite-divisibility assertion are logically distinct:
gamma mixtures of arbitrary shape larger than two need not be
infinitely divisible, as the authors point out.

\section{The principal-eigenfunction inflection question of Li, Morgan and Siegrist}
\label{sec:lms-spectral}
Y. Li, R. Morgan and M. Siegrist, \emph{Concavity structure of
fractional differential equations and zeros of Mittag--Leffler
functions}, J. Anal. (2026),
\href{https://doi.org/10.1007/s41478-026-01192-9}{doi:10.1007/s41478-026-01192-9},
specify an eigenfunction prediction immediately after their Conjecture~6.1.
For $1<a<2$ let $\lambda_a$ be the least positive $z$ such that
$E_{a,a}(-z)=0$, where $E_{a,b}(z)=\sum_{j\ge0}z^j/\Gamma(aj+b)$.
They predict that $E_{a,a-2}(-\lambda_a x^a)$ has exactly one zero
in $0<x<1$. Their Theorem~6.2 gives existence. We prove uniqueness
and simplicity. The same paper's numbered conjecture for arbitrary
nonnegative forcing has a different answer. A
\href{https://www.trureturing.org/research/li-morgan-siegrist-2026-inflection-points-conjecture-6-1/}{public counterexample}
to the latter states separately that the eigenfunction question is open.

\begin{theorem}\label{th:lms-spectral}
Let $1<a<2$, let $g(t)=t^{a-1}E_{a,a}(-t^a)$, and put
$b=\sin(\pi/a)$. For every integer $k\ge0$ and every interval
$I\subset(0,\infty)$ of length at most $\pi/b$, the derivative
$g^{(k)}$ has at most two zeros in $I$, counting multiplicities.
The principal eigenfunction $x^{a-1}E_{a,a}(-\lambda_a x^a)$
has exactly one inflection point on $(0,1)$, and the corresponding
zero of $E_{a,a-2}(-\lambda_a x^a)$ is simple.
\end{theorem}
\begin{proof}
Write $\theta=\pi/a$ and $c=-\cos\theta>0$. The transform of $g$
is $1/(1+s^a)$, by its series first on the real axis $s>1$.
Deform the inversion contour about the principal negative real cut.
The simple poles $e^{\pm i\theta}$ have residues $-e^{\pm i\theta}/a$.
For each $t>0$ this gives
\[
 g(t)=-\frac2a e^{-ct}\cos(bt+\theta)-K(t),\qquad
 K(t)=\frac{-\sin(\pi a)}{\pi}\int_0^\infty
 \frac{e^{-rt}r^a}{1+2r^a\cos(\pi a)+r^{2a}}\,dr .
\]
Indeed the cut is minus the imaginary part of the upper boundary
value of the transform, divided by $\pi$. The large arc vanishes
since $a>1$, and the small arc at zero vanishes by boundedness.
For every $k$, $K_k=(-1)^kK^{(k)}$ is the strictly positive Laplace
transform of the displayed density multiplied by $r^k$. Its logarithm
is convex by Cauchy--Schwarz or the nonnegative variance formula.

After multiplying $g^{(k)}$ by $(-1)^k$ the oscillatory term is
$s_k(t)=(-1)^{k+1}(2/a)e^{-ct}\cos(bt+(k+1)\theta)$, and the
other term is $-K_k(t)$. Zeros occur only where $s_k>0$.
Two consecutive such open lobes are separated by a gap of length
$\pi/b$, so $I$ can contain zeros from at most one lobe. Within it,
\[
 (\log(s_k/K_k))''=-b^2\sec^2(bt+(k+1)\theta)
                      -(\log K_k)''<0 .
\]
Hence there are at most two zeros counted with multiplicity: the
derivative is strictly decreasing, and a stationary zero has nonzero
second derivative. The zeros and multiplicities of $s_k-K_k$ agree
with those of this log ratio.

Let $T=\lambda_a^{1/a}$. Near zero $g(t)\sim
t^{a-1}/\Gamma(a)>0$. At $L=(3\pi/2-\theta)/b$ the cosine term
vanishes, so $g(L)=-K(L)<0$ and $T<L<\pi/b$.
On $0<x<1$ set $u(x)=T^{1-a}g(Tx)>0$. It is the principal
eigenfunction of the source, with homogeneous endpoints,
$\mathcal D_{0+}^a u=-\lambda_a u$. Its twice differentiated series
gives $u''(x)=x^{a-3}E_{a,a-2}(-\lambda_a x^a)$.
Theorem~5.2 of Li--Morgan--Siegrist applied to $F=-u$ and
$f=\lambda_a u$ gives $u''<0$ near zero and $u''>0$ near one
(their proof gives strict curvature). Its weighted derivative
hypothesis holds since $u'(x)=O(x^{a-2})$ and
$(3-a)/(2-a)+a-2>0$. The same signs follow from analyticity and
the stated strict concavity on neighborhoods, together with the
two-zero bound. The bound for $k=2$ permits at most two zeros
counting multiplicity. Opposite endpoint signs require odd total
multiplicity, hence exactly one simple zero. The positive factor
$x^{a-3}$ preserves it.
\end{proof}

\paragraph{The smallest-modulus question.}
The discussion before the proof of the source's Theorem~6.3 asks
whether the first positive zero of $E_{a,a-2}(-z)$ is a zero
of smallest modulus. In fact the unique zero of smallest modulus
is strictly negative. In $P(z)=E_{a,a-2}(z)$ the constant term
$1/\Gamma(a-2)$ is negative and every coefficient of positive
degree is strictly positive. Thus $P$ has a unique positive zero
$r$. The triangle inequality excludes all zeros in $|z|<r$.
Equality on $|z|=r$ forces the degree-one and degree-two monomials
to have the same phase, giving $z=r$. Replacing $z$ by $-z$
proves the assertion. This last observation only uses elementary
coefficient signs.

\begingroup\small
\begin{thebibliography}{99}
\setlength{\itemsep}{1pt plus 0.5pt}
\bibitem{src-gaj} K. Gajdzica, \emph{On the Bessenrodt--Ono type
inequality for a wide class of $A$-partition functions}, Ramanujan J. 66
(2025), article 80. \href{https://arxiv.org/abs/2401.16267}{2024 preprint}.
\bibitem{src-roy} A. Roy, \emph{Log-concavity and the multiplicative
properties of restricted partition functions}, Ann. Comb. (2026).
\href{https://doi.org/10.1007/s00026-026-00847-5}{Publisher's text}.
\bibitem{src-gz} D. Galvin and Y. Zhang, \emph{The domination polynomial
of powers of paths and cycles}, 2024.
\href{https://arxiv.org/abs/2408.12731}{arXiv:2408.12731}.
\bibitem{src-z} S. Zhang, \emph{Domination polynomial is unimodal for
large graphs with a universal vertex}, 2021.
\href{https://arxiv.org/abs/2111.00641}{arXiv:2111.00641}.
\bibitem{src-hsw} T. Hasebe, T. Simon and M. Wang,
\emph{Some properties of the free stable distributions},
Ann. Inst. H. Poincar\'e Probab. Statist. 56 (2020), 296--325.
\href{https://doi.org/10.1214/19-AIHP962}{Publisher's text}.
\bibitem{src-bst} A. R. Booker, A. Str\"ombergsson and H. Then,
\emph{Bounds and algorithms for the $K$-Bessel function of imaginary
order}, LMS J. Comput. Math. 16 (2013), 78--108.
\href{https://doi.org/10.1112/S1461157013000028}{Publisher's text}.
\bibitem{src-bw} \textnormal{\'A}. Baricz and S. Wu,
\emph{Sharp exponential Redheffer-type inequalities for Bessel functions},
Publ. Math. Debrecen 74 (2009), 257--278.
\href{https://doi.org/10.5486/PMD.2009.4291}{Publisher's text}.
\bibitem{src-qi} F. Qi, \emph{Completely monotonic degree of a
function involving trigamma and tetragamma functions},
AIMS Math. 5 (2020), 3391--3407.
\href{https://doi.org/10.3934/math.2020219}{Publisher's text}.
\bibitem{src-ks} M. Kwa\'snicki and T. Simon, \emph{Characterisation
of the class of bell-shaped functions}.
\href{https://arxiv.org/abs/1910.07752}{arXiv:1910.07752}.
\bibitem{src-lms} Y. Li, R. Morgan and M. Siegrist,
\emph{Concavity structure of fractional differential equations and
zeros of Mittag--Leffler functions}, J. Anal. (2026).
\href{https://doi.org/10.1007/s41478-026-01192-9}{Publisher's text}.
\end{thebibliography}\endgroup

\end{document}

%% BEGIN roy_square_verify.py
"""Exact check of a finite-core certificate for Roy's square-partition table."""
from math import lcm

core = [1, 4, 9, 16, 25, 36]
L = lcm(*core)
M = 6 * L + 20
p = [1] + [0] * M
for a in core:
    for n in range(a, M + 1):
        p[n] += p[n-a]
for h in range(4,12):
    def g(v):
        return (p[h] * p[v] if v >= 0 else 0) - p[v+h]

    for r in range(L):
        data = [g(r+k*L) for k in range(6)]
        for j in range(1,6):
            data = [y-x for x,y in zip(data,data[1:])]
            assert data[0] >= 0
    mass = sum(max(0,-g(v)) for v in range(-h,L))
    floor = min(g(v) for v in range(49-h,2*L))
    assert floor > mass and min(g(v) for v in range(L,2*L))>mass
    eq = [v for v in range(h,49) if g(v)==0 and v+h<49]
    bad = [v for v in range(h,49) if g(v)<0 and v+h<49]
    print(h,mass,floor,eq,bad)
%% END roy_square_verify.py

%% BEGIN zero_forcing_audit.cpp
#include <array>
#include <bit>
#include <cstdint>
#include <iostream>
using Bits = std::uint64_t;
constexpr unsigned N=34;
constexpr Bits ALL=(Bits(1)<<N)-1;
std::array<Bits,N> adj;
bool forces(Bits white) {
 while(white){
  Bits blue=ALL^white;
  Bits next=0;
  while(blue){
   const auto v=std::countr_zero(blue);
   blue&=blue-1;
   const Bits candidates=white & adj[v];
   if(std::has_single_bit(candidates)) next |= candidates;
  }
  if(!next) return false;
  white &= ~next;
 }
 return true;
}
int main() {
 for(unsigned v=0;v<N;++v)
  for(unsigned w=0;w<N;++w)
   if(v!=w && ((v<14)!=(w<14) ||
     (v<14 && w<14 && v/2==w/2) ||
     (v>=14 && w>=14 && (v-14)/5==(w-14)/5)))
    adj[v] |= Bits(1)<<w;
 for(int t: {4,5,6}) {
   unsigned long long count=0, total=0;
   Bits white=(Bits(1)<<t)-1;
   while(white<=ALL) {
     count+=forces(white);++total;
     Bits low=white & -white, up=white+low;
     white=up+(((white^up)/low)>>2);
   }
   std::cout <<t<<" "<<total<<" "<<count<<"\n";
 }
}
%% END zero_forcing_audit.cpp

%% BEGIN kanter_verify.py
"""Uniform compact monotonicity certificate for the Kanter phase.
Run with python-flint 0.9. All bounds use exact rational grids/rounded Arb.
"""
from math import factorial
from flint import arb as A, arb_poly as Poly, ctx
ctx.prec=176
N=384
err=A(10)**-50
# Four-term grouped coefficient majorants for n>=384, including n and radius x<=37.
# M_l=4(4l+4)(l+1)!37^(4l+3)/(4l)!; ratio <1/2 for l>=96.
assert 8*388*factorial(97)*37**387 * 10**50 < factorial(384)
assert 2 * 2 * 98 * 37**4 < 385**4
J=100; I=296
xs=[A(k)/8 for k in range(I+1)]
def alpha(j):return A(1)/5+A(j)/2000
def coeff(a,point):
    nfac=1; x0=[]; xx=[]; sgn=[]
    for n in range(1,N):
        nfac*=n
        y=n*a
        g=y.gamma()/(A.pi()*nfac)
        if point: sgn.append(g*(2*y).sin_pi()/2)
        else:
            s=y.sin_pi()
            ball=(g*(s*s)).nonnegative_part()
            x0.append(ball.upper())
            xx.append(n*ball.lower().max(0))
    if point:
        pol=Poly([1]+sgn)
        dp=pol.derivative()
        return [(pol(x),dp(x)) for x in xs]
    else:
        p0,px=Poly(x0),Poly(xx)
        return [p0(x) for x in xs],[px(x) for x in xs]
# Bounds from real-variable bilinear interpolation and the rotated contour.
# Sup of |R_xx|<=8, |R_xxx|<=16; sup |R_aa|<=2e6, |R_xaa|<=6e6.
eR=(A(8)/8**2+A(2000000)/2000**2)/8+err
eB=(A(16)/8**2+A(6000000)/2000**2)/8+err
old=coeff(alpha(0),True)
minmargin=100
for j in range(J):
    nxt=coeff(alpha(j+1),True)
    upper,lower=coeff(alpha(j).union(alpha(j+1)),False)
    for k in range(I):
        vertices=[old[k],old[k+1],nxt[k],nxt[k+1]]
        R=min(q[0].lower() for q in vertices)-eR
        B=max(q[1].upper() for q in vertices)+eB
        den=upper[k+1]+err
        # (x I'/I)*R > x R', including the x->0 limiting form.
        r=(lower[k]-err)*R/den-xs[k+1]*B.max(0)
        assert R>0 and r>0, (j,k,R,B,r)
        minmargin=min(minmargin,float(r.lower()))
    old=nxt
print("Passed",J*I,"compact rectangles; margin >",minmargin*0.999)
%% END kanter_verify.py

%% BEGIN mixture_rational.py
"""Pure rational certificate for failure of the gamma(3/2)-mixture conjecture.
python-flint 0.9; no floating point enters the check."""
from math import comb
from flint import fmpq as Q, fmpq_poly as P
def coefficient(a,n):
    N=2*n; d=-Q(2*n+1,2)
    b=[P([1])];h=[P([1])]
    for k in range(1,N+1):
        b.append((a+k)*b[-1]/(k+2))
        h.append(sum((((d+1)*j-k)*b[j]*h[k-j] for j in range(1,k+1)),P([]))/k)
    return sum(((N+1-j)*h[j] for j in range(N+1)),P([]))
def bern(p):
    d=p.degree()
    return [sum((p[i]*Q(comb(j,i),comb(d,i)) for i in range(j+1)),Q(0)) for j in range(d+1)]
for i,n in enumerate([71,55,45,35,30,29,24,23],92):
    # beta = (i+u)/100, 0<=u<=1. Every degree-2n Bernstein coefficient is negative.
    poly=coefficient(P([Q(i,100),Q(1,100)]),n)
    assert poly.degree()==2*n and max(bern(poly))<0
    print(i, n, "all Bernstein coefficients negative")
%% END mixture_rational.py

%% BEGIN trigamma_degrees.py
"""Certified positivity of the two densities in the corrected degree theorem.
Python-flint 0.9; all arithmetic exact or outward-rounded Arb."""
from math import factorial
from flint import arb as R, fmpq as Q, ctx
ctx.prec = 100
N = 80
def exact_poly():
    # b(z)=z/(1-e^-z); K=b*b-zb.  Compute K up to degree N+3.
    c = [Q.bernoulli(j)/factorial(j) if j!=1 else Q(1,2) for j in range(N+4)]
    coeff = [Q(0)]
    for m in range(N+3):
        coeff.append(sum(c[i]*c[m-i]*factorial(i)*factorial(m-i)
                         for i in range(m+1))/factorial(m+1) - c[m])
    vs, ws = [], []
    for j in range(N+1):
        ex=Q((-1)**j,factorial(j))
        exprev=-j*ex
        vs.append(coeff[j+2]*(j+2)*(j+1) - int(j==1)+2*int(j==0)-2*ex-Q(13,12)*exprev)
        ws.append(int(j==0)-(11*ex+9*exprev)/12-coeff[j+3]*(j+3)*(j+2)*(j+1))
    assert vs[:2] == ws[:2] == [0,0]
    return [list(map(R,a[2:])) for a in (vs,ws)]
P,QW=exact_poly()
def hor(v,x):
    y=v[-1]
    for d in reversed(v[:-1]): y=y*x+d
    return y
small_error=R(1500)*(R(2)/3)**(N+1) # 2000*3/4 times (2/3)^(N+1)
def small(t):
    return min(hor(P,t).lower(),hor(QW,t).lower())>small_error
H=Q(0); J=Q(0); powers=[]
for k in range(1,33):
    H+=Q(1,k*k);J+=Q(1,k**3)
    powers.append((2*(R(2).zeta()+R(H)),4*(R(3).zeta()-R(J))))
tail=R(2*100000*33**3)/R(7)**32
def medium(t):
    w=-(11+9*t)/12; v=-2-13*t/12
    b=(-t).exp(); e=R(1)
    t2=t*t; p0=t*t2/6-t2
    p1=t2/2-2*t; p2=t-2
    for k,(A,B) in enumerate(powers,1):
        p=p0+A*t-B; dp=p1+A
        v+=e*(p2-2*k*dp+k*k*p)
        w+=e*(-1+3*k*p2-3*k*k*dp+k**3*p)
        e*=b
    return min(v.lower(),w.lower())>tail
def certify(left,right,pred):
    stack=[(R(left),R(right),0)]
    accepted=0
    while stack:
        a,b,depth=stack.pop()
        if pred(a.union(b)):
            accepted+=1
            continue
        assert depth<25
        c=(a+b)/2
        stack.extend([(a,c,depth+1),(c,b,depth+1)])
    print(left,right,accepted,flush=True)
certify(0,2,small)
certify(2,32,medium)
%% END trigamma_degrees.py
